\chapter{Alterations and the normal case}
\index{alteration!envelope theorem}\index{normal variety!envelope theorem}
\index{Stein factorization!envelope theorem}
\label{chap:bj-alterations}

\section{Orientation}
\label{sec:ch42-orientation}

\subsection{What is proved here, and out of what}
\label{sec:ch42-what}

Let $X$ be a normal integral projective variety over the algebraically
closed trivially valued field $k$, of arbitrary dimension and arbitrary
characteristic, and let $\theta\in N^1(X)_{\mathbf R}$.  This chapter
proves that $\theta$ has the envelope property of
\Cref{def:bj-envelope-property}:
\[
 \begin{gathered}
 (\varphi_i)_{i\in I}\subset\PSH_X(\theta)
 \ \text{nonempty and bounded above}\\
 \Longrightarrow\quad
 \Bigl(\sup_i\varphi_i\Bigr)^*\in\PSH_X(\theta).
 \end{gathered}
\]
This is \Cref{thm:bj-normal-envelope}.  It is the fifth step of the
programme of \Cref{sec:ch37-proof-map}, and it is the last step at which the
variety is still assumed normal; the two chapters that follow remove
normality in favour of unibranchness.

The proof in this chapter adds no new regularization construction
to those of \Cref{chap:bj-smooth-characteristic}.  Its transport statement,
\Cref{thm:ch42-alteration-descent}, is an \emph{assembly} of the two
transport theorems already proved: it takes the smooth case and carries it
down to a normal target along a morphism that has to be constructed and then
dismantled.  Both operations are the substance of the chapter, and both are
more delicate than the one-line summary ``take a smooth alteration and
descend'' suggests.

The construction is de Jong's theorem.  Over a field of characteristic zero
one would take a resolution of singularities $\pi:\widetilde X\to X$,
apply \Cref{thm:bj-smooth-envelope} upstairs, and descend by
\Cref{thm:bj-birational-ep-descent}; the chapter would be four lines long
and would end there.  In positive characteristic no resolution theorem is
available in every dimension, and the substitute --- a proper surjective
generically finite morphism from a smooth variety --- need not be
birational.
Its generic degree can be any positive integer, and in positive
characteristic an arbitrary alteration may have an inseparable
function-field extension.  Although de Jong's theorem permits a generically
étale choice over the present perfect field, the argument below neither
needs nor assumes that choice.
The whole difficulty of the chapter is that this degree cannot be wished
away.

The dismantling is Stein factorization.  A proper generically finite morphism can change the function field
and have positive-dimensional exceptional fibres.  When both occur,
neither of the two descent theorems applies directly.  Stein factorization writes it as a
composite
\[
 Y\xrightarrow{\ g\ }T\xrightarrow{\ h\ }X
\]
in which the two mechanisms are cleanly separated: $g$ is projective
birational, so that \Cref{chap:bj-reflection} applies to it, and $h$ is
finite surjective between normal integral projective varieties, so that
\Cref{chap:bj-finite-descent} applies to it.  The chapter's real content is
the proof that this separation is legitimate, that the middle term $T$ has
all the properties both theorems demand --- in particular that it is
\emph{normal}, which is not part of the Stein factorization theorem and has
to be imported from the relative-normalization theorem --- and that the
numerical classes match at each stage without a stray factor of the degree.

\subsection{Why an alteration, and why not a resolution}
\label{sec:ch42-why-alteration}

The hypothesis ``arbitrary characteristic'' in
\Cref{thm:bj-final-envelope} is not decoration.  It is the reason this
chapter exists in the form it has.  It is worth being explicit about what is
and is not available.

In characteristic zero, Hironaka's theorem
(\Cref{thm:appB-resolution}) provides, for every integral variety $X$ over
a field of characteristic zero, a projective birational morphism
$\pi:\widetilde X\to X$ with $\widetilde X$ smooth.  Composing
\Cref{thm:bj-smooth-envelope} on $\widetilde X$ with
\Cref{thm:bj-birational-ep-descent} for $\pi$ proves
\Cref{thm:bj-normal-envelope} immediately, provided $X$ is normal so that
\Cref{chap:bj-reflection} applies.  This is recorded as
\Cref{cor:ch42-char-zero-shortcut} below, and a reader interested only in
characteristic zero may stop there.

In positive characteristic no resolution theorem is available in every
dimension.  No argument for the arbitrary-dimensional statement of this
book may therefore assume it.  What \emph{is} available in every
characteristic and every dimension is de Jong's theorem: there is a proper
surjective generically finite morphism onto $X$ whose source is smooth.
The cost permitted by \Cref{def:bj-alteration} is a finite extension
of the function field, which need not have degree one.

Three consequences of that cost organize the chapter.

\emph{The common function field is gone.}  Every statement in
\Cref{chap:bj-reflection} rests on the identification
$K(Y)=K(X)$ and on the resulting identification of the divisorial loci
$X^{\divis}$ and $Y^{\divis}$, through
\Cref{lem:ch39-beta-an}\,\cref{itm:ch39-ban:val}.  For an alteration of
degree $d>1$ this identification does not exist, so the birational reflection theorem does not apply to that morphism.

\emph{The fibres need not be finite.}  The finite descent proof in
\Cref{chap:bj-finite-descent} first flattens its finite map and then
uses \eqref{eq:ch41-average-shape} for the resulting constant-rank
finite flat morphism.  A nonfinite alteration may
contract subvarieties, so some of its fibres can be positive-dimensional,
and the sum that defines the average then has infinitely many terms.

\emph{Neither defect is repaired by the other.}  An alteration is not a
birational morphism composed with nothing, nor a finite morphism composed
with nothing; it is genuinely both at once, and the point of Stein
factorization is that ``both at once'' can be written as ``one after the
other''.  \Cref{ex:ch42-no-separation} carries all three of these
observations to explicit functions on an explicit surface.

\subsection{Why the two mechanisms must be separated, and in which order}
\label{sec:ch42-why-separate}

Suppose that $f:Y\to X$ is neither birational nor finite.
The two descent theorems available here each have a hypothesis it fails.

The birational mechanism of \Cref{chap:bj-reflection} uses Cartier
approximants and homogeneous reflection over the common function field,
followed by carrier reflection for general psh functions.
It does not identify every homogeneous psh function with a Cartier
b-divisor; \Cref{rem:ch39-cartier-not-enough} records that distinction.
For a nontrivial field extension the common divisorial locus is absent.
The finite mechanism of
\Cref{chap:bj-finite-descent} moves a psh function down by averaging over
the fibres with the weights \eqref{eq:ch41-weights}; the averaging operator
is constructed from the determinant of a finite locally free algebra, and
there is no such algebra when a fibre is a curve.

Stein factorization repairs both defects at once, and it does so in a
specific order which is not a matter of taste.  The factorization
$f=h\circ g$ produced by \Cref{thm:appB-stein-factorization} always has
the connected-fibre factor \emph{on top} and the finite factor
\emph{underneath}.  Consequently the descent must run
\[
 \operatorname{EP}\ \text{on}\ Y
 \ \xrightarrow{\ \text{birational}\ }\
 \operatorname{EP}\ \text{on}\ T
 \ \xrightarrow{\ \text{finite}\ }\
 \operatorname{EP}\ \text{on}\ X,
\]
and never in the other order.  This is not because the other order would be
logically invalid --- both descent theorems are implications from source to
target, and either could in principle be applied first --- but because the
other order requires a factorization of $f$ with the finite map on top,
and Stein factorization does not supply that reverse-order
factorization for the given source.  \Cref{ex:ch42-second-factorization} exhibits a
special position in which such a factorization happens to exist, and shows
that even then it produces a different middle variety; the canonical one is
the Stein middle term.

\subsection{The firewall}
\label{sec:ch42-firewall}

This chapter is the fifth input of \Cref{thm:bj-final-envelope}, and the
circularity it must avoid is the shortest of all those in
\Crefrange{chap:bj-cartier-realization}{chap:bj-subintegral-normalization}:
a normal variety is unibranch, so a single appeal to
\Cref{thm:bj-final-envelope} would prove \Cref{thm:bj-normal-envelope} in
one line, and would prove nothing at all.

\begin{warning}{The envelope theorem is not available in this chapter}
\label{warn:ch42-firewall}
Nothing in \Cref{chap:bj-alterations} may use, directly or transitively:
\begin{enumerate}
\item \Cref{thm:bj-final-envelope}, in any of its formulations, nor
      \Cref{def:bj-envelope-property-final}, nor
      \Cref{cor:bj-historical-conjecture};
\item the envelope property $\operatorname{EP}$ of
      \Cref{def:bj-envelope-property} as a \emph{theorem} for any class on
      any variety, with the single exception of
      \Cref{thm:bj-smooth-envelope} on a \emph{smooth} source; wherever
      $\operatorname{EP}$ occurs as an assumption it is displayed in the
      statement carrying it, and it is always an assumption on the
      \emph{source} of a morphism;
\item closure of any class $\PSH_X(\theta)$ under suprema, regularized or
      not;
\item the obstacle envelope $P_\theta$ of
      \Cref{def:bj-obstacle-envelope} regarded as an operator with values in
      a psh class;
\item unconditional birational reflection, that is,
      \Cref{thm:bj-carrier-reflection} with the carrier hypothesis deleted,
      and \eqref{eq:bj-ordinary-reflection-under-ep} for a variety whose
      envelope property has not been separately established;
\item any statement of \Cref{chap:bj-universal-homeomorphisms},
      \Cref{chap:bj-subintegral-normalization}, or
      \Cref{chap:bj-envelope-consequences}.
\end{enumerate}
The unibranchness of a normal variety, which is true and elementary, may
therefore never be combined with \Cref{thm:bj-final-envelope} in this
chapter.
\end{warning}

\begin{remark}[Why the firewall is not vacuous here]
\label{rem:ch42-firewall-not-vacuous}
Three places in this chapter are tempting.

First, \Cref{thm:bj-normal-envelope} would follow from
\Cref{thm:bj-final-envelope} in one line, because a normal variety is
unibranch: its local rings are normal domains, hence its normalization is
the identity and has one point over each point.  This is the shortest
possible circle in the book, and it is why the firewall is stated before
anything is proved.

Second, the proof of \Cref{thm:ch42-alteration-descent} applies
\Cref{thm:bj-birational-ep-descent} to $g:Y\to T$ and needs
$\operatorname{EP}_Y(g^*h^*\theta)$.  It would be comfortable to say that
$Y$ is normal projective and quote the normal case --- which is the
theorem being proved.  Instead the hypothesis is supplied by
\Cref{thm:bj-smooth-envelope}, whose proof
(\Cref{chap:bj-smooth-characteristic}) is independent of this chapter and
uses no envelope property at all.

Third, \Cref{cor:ch42-birational-bijection} and
\Cref{cor:ch42-finite-bijection} state the pullback identifications for normal projective varieties:
a bijection along a projective birational morphism, and a bijection onto
the fibrewise-constant subclass along a finite surjective morphism.  These are consequences of \Cref{thm:bj-normal-envelope} and are
stated only after it; they are not used in its proof, and in particular
\eqref{eq:bj-ordinary-reflection-under-ep} is never invoked before its
hypothesis has been established for the variety in question.
\end{remark}

\subsection{The five tasks, and where each is discharged}
\label{sec:ch42-tasks}

\begin{enumerate}
\item\label{itm:ch42-task:alteration}
      \emph{Alterations in arbitrary characteristic.}  State the exact
      form of de Jong's theorem that is consumed here, with its hypotheses
      displayed and with the passage from ``regular'' to ``smooth'' and
      from ``quasi-projective'' to ``projective'' carried out rather than
      asserted.  This is \Cref{thm:bj-de-jong-input}, in
      \Cref{sec:ch42-de-jong}, preceded by the degree calculus of
      \Cref{sec:ch42-generic-degree}.
\item\label{itm:ch42-task:stein}
      \emph{Stein factorization.}  Produce the middle term, prove its
      universal property, and prove that it is a normal integral projective
      variety.  Normality is the item that is \emph{not} part of the Stein
      theorem and must be imported separately.  This is
      \Cref{lem:ch42-stein-universal} and \Cref{lem:bj-alteration-middle},
      in \Cref{sec:ch42-stein}.
\item\label{itm:ch42-task:separate}
      \emph{Separating the birational and finite steps.}  Prove, and not
      merely assert, that $g$ is birational and $h$ finite surjective,
      that both are morphisms of normal integral projective varieties, that
      the separation is canonical, and that each factor is trivial exactly
      when the corresponding feature of $f$ is absent.  This is
      \Cref{prop:ch42-separation}, in \Cref{sec:ch42-separation}.
\item\label{itm:ch42-task:descent}
      \emph{Descent from a smooth alteration.}  Prove the general descent
      statement: if $f:Y\to X$ is a proper surjective generically finite
      morphism from a normal integral projective variety onto a normal
      integral projective variety, then $\operatorname{EP}_Y(f^*\theta)$
      implies $\operatorname{EP}_X(\theta)$.  This is
      \Cref{thm:ch42-alteration-descent}, in \Cref{sec:ch42-descent}, and it
      assembles the two preceding descent implications.
\item\label{itm:ch42-task:normal}
      \emph{The envelope property for normal projective varieties.}
      Combine the previous item with \Cref{thm:bj-smooth-envelope} and de
      Jong's theorem.  This is \Cref{thm:bj-normal-envelope}, in
      \Cref{sec:ch42-normal-theorem}, together with its consequences for
      normalizations, towers, and quotients.
\end{enumerate}

\subsection{Standing hypotheses and notation}
\label{sec:ch42-hypotheses}

Throughout the chapter, $k$ is an algebraically closed field equipped with
the trivial absolute value, of arbitrary characteristic.  All varieties are
integral and projective over $k$ unless the contrary is said; ``variety''
carries the meaning fixed in the conventions, namely a separated integral
scheme of finite type over $k$.  The dimension is arbitrary.  Numerical
classes are real: $\theta\in N^1(X)_{\mathbf R}$, with no positivity
hypothesis whatever.  Envelope-property statements for an empty psh class are vacuous,
and \Cref{lem:ch34-psh-nonempty} shows that nonemptiness requires
pseudoeffectivity.  The geometric lemmas are independent of this condition.
Nefness is needed when existence of a function bounded on both sides
is asserted; the ample surface example below supplies such a carrier.

The psh class $\PSH_X(\theta)$ is the Boucksom--Jonsson class of
\Cref{def:bj-psh}: upper semicontinuous, generically finite, and a pointwise
limit of a \emph{decreasing} net of generically finite Fubini--Study
functions whose classes converge to $\theta$ but need not equal it.  The
permission for the classes to vary is used at every application of
\Cref{cor:bj-coefficient-norm-limit} inside
\Cref{chap:bj-finite-descent}, and it is what makes the chapter's freedom
from positivity hypotheses possible; see
\Cref{rem:ch41-varying-classes-needed}.

We write $\operatorname{EP}_X(\theta)$ for the envelope property of
\Cref{def:bj-envelope-property}, $X^{\An}$ for the Berkovich
analytification with respect to the trivial absolute value, $X^{\divis}$
for the trivial valuation together with the positive real multiples of
divisorial valuations of $K(X)$, and $v_{X,\triv}$ for the trivial
valuation of $K(X)$, viewed as a point of $X^{\An}$.  For a semivaluation
$x\in X^{\An}$ we write $s(x)$ for its support and $c(x)$ for its
centre.

Finally, a word on a notational trap.  The letter $d$ always denotes the
generic degree $[K(Y):K(X)]$ of the alteration, equivalently the degree of
its finite Stein factor; it never denotes a dimension.  The letter $g$
always denotes the birational Stein factor and never a genus.

\subsection{What is used from earlier chapters}
\label{sec:ch42-dependencies}

From \Cref{chap:bj-fs-pl-psh} we use \Cref{def:bj-psh}, the existence
threshold \Cref{lem:ch34-psh-nonempty}, generic finiteness
\Cref{prop:ch34-generic-finiteness}, and the elementary calculus of
$\log|s|$ in \Cref{lem:ch34-log-s-properties}.  From
\Cref{chap:bj-envelope-history} we use the definition
\Cref{def:bj-envelope-property}.  The b-divisor comparison,
regularization, and carrier arguments are not reopened here; they enter only
through the two descent theorems listed next.

The three inputs proper are:
\begin{itemize}
\item \Cref{thm:bj-smooth-envelope}: the envelope property for a smooth
      projective variety and an arbitrary real class, in arbitrary
      characteristic;
\item \Cref{thm:bj-birational-ep-descent}: birational descent of the
      envelope property along a projective birational morphism of normal
      integral projective varieties, together with the tower form
      \Cref{cor:ch39-ep-tower};
\item \Cref{thm:bj-finite-normal-descent}: finite descent along a finite
      surjective morphism of normal integral projective varieties, together
      with \Cref{cor:ch41-average-pullback},
      \Cref{cor:ch41-real-weights} and
      \Cref{lem:ch41-proper-gf-pullback}, all of which are used through it
      rather than directly.
\end{itemize}
None of the three uses any statement of the present chapter, and none uses
\Cref{thm:bj-final-envelope}; the three firewalls
\Cref{warn:ch39-firewall}, \Cref{warn:ch40-firewall} and
\Cref{warn:ch41-firewall} record this.

The geometric inputs are collected in \Cref{app:birational-geometry}:
Stein factorization \Cref{thm:appB-stein-factorization} with its normality
and birationality complements \Cref{prop:appB-stein-normal} and
\Cref{lem:appB-stein-birational}, and de Jong's theorem
\Cref{thm:appB-de-jong}.  The exact Stacks Project statements behind them,
the hypotheses to be checked at each use, and the distinction between direct
and inherited dependencies are audited in \Cref{app:stacks}; no tag is
quoted here that is not listed there.  The flattening theorem
\Cref{thm:appB-flattening} is used only inside
\Cref{chap:bj-finite-descent} and does not reappear here, a point worth
keeping in view because the flattening square of
\eqref{eq:bj-flatification-square} and the Stein factorization of this
chapter are different diagrams built for different purposes
(\Cref{warn:ch42-stein-not-flattening}).

\section{Alterations: the precise input}
\label{sec:ch42-alterations}

The purpose of this section is to fix exactly what is meant by an alteration,
to establish the elementary degree calculus that the rest of the chapter
uses without comment, and to state de Jong's theorem in the form in which it
is consumed --- that is, with the passage from a regular quasi-projective
source to a smooth projective one carried out, rather than absorbed into the
phrase ``choose a smooth projective alteration''.

\subsection{Generic finiteness and degree}
\label{sec:ch42-generic-degree}

\begin{definition}[Alteration]
\label{def:bj-alteration}
An \emph{alteration} of an integral variety $X$ is a morphism
\[
                         f:Y\longrightarrow X
\tag{42.1}\label{eq:bj-alteration-map}
\]
that is proper, surjective, and generically finite, with $Y$ integral.
It need not be birational.
Its generic degree may be greater than one and may have an inseparable part.
\end{definition}

The distinction between an alteration and a resolution is essential.
A resolution is birational, whereas an alteration is allowed to change the
function field by a finite extension.
This flexibility is what makes alterations available in arbitrary
characteristic.
It also explains why one cannot apply birational reflection directly to
\eqref{eq:bj-alteration-map}.

Before anything else, the phrase ``generically finite'' has to be pinned
down, because three inequivalent-looking formulations of it occur in the
literature and the proofs below use all three.  They are in fact equivalent
for a dominant morphism of integral varieties, and it is worth proving this
once.

\begin{lemma}
\label{lem:ch42-generic-finite}
Let $f:Y\to X$ be a dominant morphism of integral $k$-varieties, and let
$\eta$ be the generic point of $X$.  The following are equivalent.
\begin{enumerate}
\item\label{itm:ch42-gf:dim} $\dim Y=\dim X$.
\item\label{itm:ch42-gf:field} The field extension
      $K(X)\subseteq K(Y)$ induced by $f$ is finite.
\item\label{itm:ch42-gf:fibre} The generic fibre $Y_\eta$ is finite over
      $K(X)$; in fact $Y_\eta=\Spec K(Y)$.
\end{enumerate}
If moreover $f$ is proper, these conditions are equivalent to:
\begin{enumerate}
\setcounter{enumi}{3}
\item\label{itm:ch42-gf:open} There is a nonempty open $U\subseteq X$ such
      that $f^{-1}(U)\to U$ is finite and surjective.
\end{enumerate}
One may take $U=X\setminus f(Y\setminus Y_{\mathrm{qf}})$,
where $Y_{\mathrm{qf}}$ is the quasi-finite locus in the source;
this is the largest target open over which the entire restriction
of $f$ is quasi-finite.
\end{lemma}

\begin{proof}
Since $f$ is dominant and $Y$ is integral, the generic point of $Y$
lies over $\eta$, so $Y_\eta$ is nonempty; it is the localization of
$Y$ at the multiplicative system of nonzero pullbacks from $X$, hence is
an integral scheme of finite type over $K(X)$ with function field
$K(Y)$.  The dimension formula for a finitely generated field extension
gives
\[
 \dim Y=\operatorname{trdeg}_kK(Y)
       =\operatorname{trdeg}_{K(X)}K(Y)+\operatorname{trdeg}_kK(X)
       =\dim Y_\eta+\dim X .
\]
Hence \cref{itm:ch42-gf:dim} holds if and only if $\dim Y_\eta=0$.

\cref{itm:ch42-gf:dim}$\Rightarrow$\cref{itm:ch42-gf:fibre}: an integral
scheme of finite type and dimension zero over a field $F$ is the spectrum
of a finite field extension of $F$; applied to $Y_\eta$, this gives
$Y_\eta=\Spec K(Y)$ with $K(Y)/K(X)$ finite.

\cref{itm:ch42-gf:fibre}$\Rightarrow$\cref{itm:ch42-gf:field} is contained in
the previous sentence, and
\cref{itm:ch42-gf:field}$\Rightarrow$\cref{itm:ch42-gf:dim} follows from the
displayed dimension formula, a finite extension having transcendence degree
zero.

Suppose now that $f$ is proper.  Under the first three equivalent
conditions the quasi-finite locus in $Y$ contains the generic point.  Its
complement is closed, and properness makes its image closed in $X$; removing
that image gives a nonempty open $U$ over which $f$ is quasi-finite.
The restriction $f^{-1}(U)\to U$ is proper and quasi-finite, hence finite,
and it is surjective because $f$ is dominant and proper.  This proves
\cref{itm:ch42-gf:open}.  Conversely, the generic fibre of a finite
morphism is finite, so \cref{itm:ch42-gf:open} implies
\cref{itm:ch42-gf:fibre}.
\end{proof}

\begin{definition}[Degree of an alteration]
\label{def:ch42-alteration-degree}
The \emph{degree} of an alteration $f:Y\to X$ is
\[
 \deg f:=[K(Y):K(X)]
       =[K(Y):K(X)]_{\mathrm s}\cdot[K(Y):K(X)]_{\mathrm i},
\]
the product of the separable and the inseparable degrees.  The alteration is
\emph{birational} if $\deg f=1$, \emph{separable} if
$[K(Y):K(X)]_{\mathrm i}=1$, and \emph{purely inseparable} if
$[K(Y):K(X)]_{\mathrm s}=1$.  In characteristic zero every alteration is
separable.
\end{definition}

This decomposition is not idle bookkeeping, but it must not be read as a
dictionary.  The separable degree does \emph{not} count the points of a
fibre, and the weights $e(w/z)$ of \eqref{eq:ch41-weights} are not the
inseparable degree: by \Cref{rem:ch41-weight-anatomy} each weight collects
a local Artin length and the full degree of a completed valued-field
extension, including ramification, residue degree and possible defect.
The product of just the first two field invariants is asserted
only under the defectless hypothesis.  Already \Cref{ex:ch42-endtoend}, in characteristic
$\ne2$, has $d=2$ with $d_{\mathrm s}=2$ and $d_{\mathrm i}=1$,
and yet a single point above the divisorial point $v_P$, carrying the
tame weight $e=2$.  What survives, and all that survives, is
$\sum_{w\mid z}e(w/z)=d$.  The example that forces the bookkeeping in
characteristic $p$ is \Cref{ex:ch42-inseparable-cone}.

\begin{lemma}
\label{lem:ch42-degree-composite}
Let $Z\xrightarrow{\ b\ }Y\xrightarrow{\ a\ }X$ be dominant morphisms of
integral $k$-varieties.
\begin{enumerate}
\item\label{itm:ch42-deg:comp} If $a$ and $b$ are generically finite,
      so is $a\circ b$, and
      $\deg(a\circ b)=\deg a\cdot\deg b$.
\item\label{itm:ch42-deg:cancel} If $a\circ b$ is generically finite, then
      so are $a$ and $b$.
\item\label{itm:ch42-deg:proper} If $a\circ b$ is proper and $a$ is
      separated, then $b$ is proper.
\item\label{itm:ch42-deg:alteration} If $a\circ b$ is an alteration and
      $a$ is separated and dominant with $Y$ integral, then $b$ is an
      alteration of $Y$ and $a$ is an alteration of $X$.
\end{enumerate}
\end{lemma}

\begin{proof}
\cref{itm:ch42-deg:comp} and \cref{itm:ch42-deg:cancel} are the multiplicativity
of degrees in a tower $K(X)\subseteq K(Y)\subseteq K(Z)$ of field
extensions, together with
\Cref{lem:ch42-generic-finite}\,\cref{itm:ch42-gf:field}: the tower is finite
if and only if both steps are, and then the degrees multiply.

\cref{itm:ch42-deg:proper} is the standard cancellation property of proper
morphisms: $b$ is the composite of its graph
$\Gamma_b:Z\to Z\times_XY$, a closed immersion because $a$ is separated,
with the second projection $Z\times_XY\to Y$, which is proper as a base
change of the proper morphism $a\circ b$.

\cref{itm:ch42-deg:alteration}  By \cref{itm:ch42-deg:proper}, $b$ is proper;
by \cref{itm:ch42-deg:cancel}, both $a$ and $b$ are generically finite;
$b$ is surjective because its image is closed and $b$ is dominant
by hypothesis.  Then $a$ is surjective since $a\circ b$ is.
For universal closedness, make any base change $X'\to X$ and take
a closed subset $C\subset Y_{X'}$.  Its inverse image in $Z_{X'}$
is closed, and its image under the proper composite is exactly
$a_{X'}(C)$, since the base change of the proper surjection $b$
is surjective.  Thus $a$ is universally closed.  It is separated
and of finite type, hence proper.
\end{proof}

\begin{example}[Frobenius is an alteration and is not a resolution]
\label{ex:ch42-frobenius-alteration}
Let $\operatorname{char}k=p>0$ and let $X$ be an integral projective
$k$-variety of dimension $n$.  Its relative Frobenius
$F_{X/k}:X\to X^{(p)}$ is a finite surjective $k$-morphism and a
universal homeomorphism.  (The absolute Frobenius is not in general a
$k$-morphism, and no identification $X^{(p)}\simeq X$ is needed.)  Since
$k$ is perfect, the induced extension of function fields has degree
\[
 [K(X):F_{X/k}^*K(X^{(p)})]=p^{\,n},
\]
because $K(X)$ is a finitely generated field of transcendence degree $n$
over the perfect field $k$, hence has $p$-degree $n$.  Thus
$F_{X/k}$ is an alteration of degree $p^n$, purely inseparable, and it
is birational only when $n=0$.

The example is worth keeping in mind for two reasons.  First, it shows that
an alteration can be finite, so that the birational factor of its Stein
factorization is an isomorphism; the two factors of the factorization
genuinely are independent of each other.  Second, it shows that ``the
alteration is close to the identity'' is a false intuition:
$F_{X/k}$ is a universal homeomorphism and yet has generic degree
$p^n$.  The field degree $p^n$ must remain in the algebraic bookkeeping.
If $X$ is smooth, relative Frobenius is finite locally free of
that rank, and its unique analytic lift has weight $p^n$
in \Cref{lem:ch41-fibre-decomposition}; see
\Cref{ex:ch41-frobenius-average,ex:appB-frobenius-alteration}.
For singular $X$, the finite morphism need not be flat.
The constant-rank weights of Chapter~41 then belong to its flattened
map, not to dimensions of the original scheme fibres.
\end{example}

\subsection{De Jong's theorem, with its hypotheses displayed}
\label{sec:ch42-de-jong}

We use de Jong's theorem in the following form.

\begin{theorem}\label{thm:bj-de-jong-input}
Let $X$ be an integral variety over a field $k$, and let
$Z\subsetneq X$ be closed.
There are an alteration $f:Y\to X$ and an open immersion
$j:Y\hookrightarrow\overline Y$ such that $\overline Y$ is regular and
projective and
\[
       \overline{f^{-1}(Z)}\cup(\overline Y\setminus Y)
\]
is the support of a strict normal-crossings divisor.
The construction may pass through a finite extension of the ground field;
when the original field is perfect, the regular projective source is smooth
over the relevant field, and the alteration may be chosen generically
étale.
\end{theorem}

\begin{proof}[Imported foundational core and specialization]
The existence assertion, including the prescribed strict-normal-crossings
boundary, is exactly Theorem~4.1 of \cite{deJong1996}; this is the imported
foundational core.  De Jong's Remark~4.2 gives the more precise field
statement: the construction factors through a finite extension
$k_1/k$, the projective compactification is geometrically integral and
smooth over $k_1$, and, when $k$ is perfect, it is smooth over $k$.
These anchors prove every assertion in the stated general form.
The final sentence of de Jong's Theorem~4.1 supplies the generically étale
choice when $k$ is perfect.

Only the case $Z=\varnothing$ is needed below.
Because $k$ is algebraically closed, it has no nontrivial finite algebraic
extension.
Because it is perfect, regularity of a finite-type $k$-scheme is equivalent
to smoothness over $k$.

There is one small properness point hidden by the customary phrase
``choose a smooth projective alteration.''
The alteration $Y\to X$ is proper, and $X\to\Spec k$ is proper, so $Y$
is proper over $k$.
The open immersion $j:Y\hookrightarrow\overline Y$ is therefore proper,
since $\overline Y$ is separated over $k$.
Its image is both open and closed.
As $\overline Y$ is integral and $Y$ is nonempty, the image is all of
$\overline Y$.
Thus, in the present projective situation, de Jong's output is indeed a
smooth projective alteration
\[
                         f:Y\longrightarrow X.
\tag{42.2}\label{eq:bj-smooth-projective-alteration}
\]
\end{proof}

Four points in this statement repay a closer look, because each of them is a
place where a hypothesis could be, and in the literature sometimes is,
silently dropped.

\emph{Regular is not smooth.}  Over a field that is not perfect, a regular
finite-type scheme need not be smooth: the standard example is
$\Spec\bigl(\mathbf F_p(t)[x]/(x^p-t)\bigr)$, the spectrum of a field,
hence regular, but not smooth over $\mathbf F_p(t)$ because its module of
differentials is one-dimensional over that field while the scheme has
dimension zero.  De Jong's construction produces a regular source; the
passage to smoothness is a separate step, and it is legitimate here only
because $k$ is algebraically closed, hence perfect.

\emph{The finite ground-field extension is not harmless in general.}
De Jong's construction may replace $k$ by a finite extension $k_1$, and
the resulting $Y$ is then a $k_1$-variety.  For the trivially valued
theory of \Cref{part:bj} this would be a genuine change of setting: the
analytification, the scaling action, and the class $\PSH_X(\theta)$ are all
defined relative to the ground field, and a change of ground field would
require a comparison theorem that is not proved in this book.  The standing
hypothesis that $k$ is algebraically closed removes the ground-field
issue without requiring a new comparison theorem.

\emph{Projectivity of the source is not automatic from the statement.}
De Jong's theorem provides an open immersion of $Y$ into a regular
projective $\overline Y$; it does not say that $Y$ itself is projective.
The argument given in the proof --- $Y$ is proper over $k$, hence the
open immersion is proper, hence its image is closed as well as open, hence
everything --- is where the projectivity of $X$ is used.  For a nonproper $X$, the source need not be proper over $k$ and can
remain an open part of its projective compactification.
The present BJ argument requires projective varieties.

\emph{Projectivity of the morphism.}  We also record that the morphism
$f$ itself is projective, not merely proper: its source is projective over
$k$, its target is separated, and a proper morphism from a projective
$k$-scheme to a separated $k$-scheme is projective, by taking the closed
graph in $\mathbf P^N\times X$.  This is used in
\Cref{lem:bj-alteration-middle} to see that the birational Stein factor is
projective, which is what \Cref{thm:bj-birational-ep-descent} requires.

\begin{warning}{What is neither controlled nor used}
\label{warn:ch42-de-jong-scope}
\Cref{thm:bj-de-jong-input} gives a smooth projective source and a proper
surjective generically finite morphism.  The argument never requires:
\begin{enumerate}
\item any bound on the degree $d=\deg f$; in particular $d$ is not
      controlled by the singularities of $X$, and no statement below is
      allowed to depend on $d$ being small, prime, or invertible in $k$;
\item the generically étale choice: over the present perfect field the
      theorem permits a generically étale choice, but the descent theorem
      is proved for an arbitrary alteration and therefore also covers
      inseparable function-field extensions such as
      \Cref{ex:ch42-inseparable-cone};
\item any relation between the exceptional locus of $f$ and the singular
      locus of $X$;
\item a canonical or functorial choice of alteration; no compatibility
      between different choices is used.
\end{enumerate}
In particular the strict normal-crossings boundary in
\Cref{thm:bj-de-jong-input}, which is the substantive content of de Jong's
theorem and the reason its proof is hard, is \emph{not} used in this
chapter.  What is used is only the existence of a smooth integral projective
source.
\end{warning}

\begin{remark}[What is and is not used]
\label{rem:ch42-what-is-used}
The normal-crossings boundary in de Jong's theorem is not used here.
Nor do we use the generically étale choice that is available over $k$.
The proof needs only a smooth integral projective source and a proper,
surjective, generically finite map to $X$.
Consequently no separability hypothesis is smuggled into the descent.
\end{remark}

\begin{remark}[The status of resolution in positive characteristic]
\label{rem:ch42-resolution-status}
It is natural to ask why the weaker statement is used when the stronger one
would shorten the chapter.  The answer is that the stronger statement is not
a theorem.  Resolution of singularities in characteristic zero is Hironaka's
theorem \cite{Hironaka1964I,Hironaka1964II}, imported as
\Cref{thm:appB-resolution} and used in this book only over fields of
characteristic zero.  In positive characteristic no resolution theorem is
available in every dimension.  Since \Cref{thm:bj-final-envelope} claims
arbitrary dimension and arbitrary characteristic, no step of its proof may
use resolution, and the present chapter is precisely the place where the
all-dimensional substitute is installed.

The substitution is not free.  A resolution would give a birational morphism
and the chapter would consist of a single application of
\Cref{thm:bj-birational-ep-descent}; see
\Cref{cor:ch42-char-zero-shortcut}.  What the substitution costs is the
whole of \Cref{chap:bj-finite-descent}: determinant norms, flattening,
weights that are Artin lengths rather than counts of points, and the
analytic openness of a finite morphism of normal varieties.  In the ledger
of the book, \Cref{chap:bj-finite-descent} is the price of doing without
Hironaka.
\end{remark}

\subsection{The degree cannot be ignored}
\label{sec:ch42-not-resolution}

\begin{example}[Why alteration degree cannot be ignored]
\label{ex:ch42-degree-matters}
Let $X$ already be smooth and let $f:Y\to X$ be a finite cover of degree
$d>1$.
The map is an alteration but is not birational.
The common function-field identification used in birational reflection does
not exist.
The finite descent theorem instead records the degree through determinant
norms and the normalized fibre average.
This simple example is the reason for retaining the finite factor in the
Stein decomposition below.
\end{example}

The following example makes the same point in a form that can be checked
against the definitions: it exhibits, on an explicit smooth alteration, a
psh function upstairs that is visibly not the pullback of anything
downstairs, so that no descent statement can be proved by pretending the
degree away.

\begin{example}[A psh function upstairs that is not a pullback]
\label{ex:ch42-not-a-pullback}
Let $\operatorname{char}k\ne2$, let $X=\mathbf P(1,1,2)$ be the quadric
cone of \Cref{ex:ch41-nonflat-cone}, with weighted coordinates $a,b$ of
weight one and $z$ of weight two, and let
\[
 h:T=\mathbf P^2\longrightarrow X,\qquad [u:v:w]\longmapsto[u:v:w^2]
\]
be the degree-two quotient by $\iota[u:v:w]=[u:v:-w]$.  Put
$L:=\cO_{\mathbf P^2}(1)$ and
\[
 \varphi:=\log|u-w|\in\cH^{\mathrm{gf}}(L)\subset\PSH_T(c_1(L)).
\]
Then $\varphi$ is not of the form $h^*\psi$ for any function $\psi$ on
$X^{\An}$.  Indeed, $\iota^*(u-w)=u+w$, so $\varphi\circ\iota^{\An}
=\log|u+w|\ne\varphi$; on the other hand every pullback $h^*\psi$ is
$\iota$-invariant, because $h\circ\iota=h$.  Concretely, write
$p_\pm:=[1:0:\pm1]\in\mathbf P^2$, two distinct closed points because
$\operatorname{char}k\ne2$, and let $y\in X^{\An}$ be the trivial
semivaluation $v_{q,\triv}$ at the closed point $q:=h(p_+)=h(p_-)
=[1:0:1]\in X$.  Its two preimages under $h^{\An}$ are
$v_{p_+,\triv}$ and $v_{p_-,\triv}$, and
\[
 \varphi\bigl(v_{p_+,\triv}\bigr)=-\infty,
 \qquad
 \varphi\bigl(v_{p_-,\triv}\bigr)=0,
\]
the first because $u-w$ vanishes at $[1:0:1]$, the second because it does
not vanish at $[1:0:-1]$ and the valuation is trivial.  A pullback takes
the same value at both.

The example is the reason for the whole of \Cref{chap:bj-finite-descent}:
descending along a morphism of degree $d>1$ is not a matter of recognizing
which functions upstairs came from downstairs, because a general section presentation need not descend.  It is a matter of \emph{producing} a function
downstairs, and the producer is the determinant norm.  The average of the
present $\varphi$ is computed in \Cref{ex:ch41-descent-cone}: it is
$\tfrac12\log|a^2-z|$, a Fubini--Study function on $X$ of class
$\tfrac12c_1(H)$, and it is not $\varphi$ transported by any bijection.
\end{example}

\section{Stein factorization of an alteration}
\label{sec:ch42-stein}

Let $X$ be normal and let $f:Y\to X$ be the smooth projective alteration
in \eqref{eq:bj-smooth-projective-alteration}.
The Noetherian Stein factorization gives
\[
                 Y\xrightarrow{g}T\xrightarrow{h}X,
                 \qquad
                 T=\Spec_X(f_*\cO_Y),
\tag{42.3}\label{eq:bj-alteration-stein}
\]
with $g$ proper with geometrically connected fibres, $h$ finite, and
$g_*\cO_Y=\cO_T$.
The statement, including the identification of $T$ with the normalization
of $X$ in $Y$, is
\cite[\href{https://stacks.math.columbia.edu/tag/03H0}{Tag 03H0}]
{stacks-project}, imported as \Cref{thm:appB-stein-factorization} and
audited in \Cref{app:stacks}.

The phrase ``normalization of $X$ in $Y$'' is relative.
It does not say that $T=X^\nu$, the normalization of $X$ in $K(X)$.
Rather, the relevant function field will be $K(Y)$, and the finite map
$T\to X$ generally has degree $[K(Y):K(X)]$.

\subsection{The universal property, and why it matters}
\label{sec:ch42-stein-universal}

Everything this chapter needs about $T$ follows from one formal statement:
$T$ is initial among the finite --- indeed among the affine --- factors of
$f$.  We prove it, because the canonicity of the factorization is what
makes the phrase ``the birational part'' and ``the finite part'' of an
alteration well defined, and because \Cref{ex:ch42-second-factorization}
exhibits a competing factorization that is not the Stein one.

\begin{lemma}
\label{lem:ch42-stein-universal}
Let $f:Y\to X$ be a proper morphism with $X$ locally noetherian, and let
$f=h\circ g$ be its Stein factorization \eqref{eq:bj-alteration-stein}.
\begin{enumerate}
\item\label{itm:ch42-su:initial}
      For every factorization $f=h'\circ g'$ with
      $h':X'\to X$ affine, there is a unique $X$-morphism
      $c:T\to X'$ with $g'=c\circ g$.
\item\label{itm:ch42-su:unique}
      If in addition $h'$ is finite and $g'_*\cO_Y=\cO_{X'}$, then
      $c$ is an isomorphism.
\item\label{itm:ch42-su:finite}
      $g$ is an isomorphism if and only if $f$ is affine, and then
      $f=h$; in particular this holds when $f$ is finite.
\end{enumerate}
\end{lemma}

\begin{proof}
\cref{itm:ch42-su:initial}  Since $h'$ is affine, $X'=\Spec_X\mathcal
A'$ with $\mathcal A':=h'_*\cO_{X'}$, and for any $X$-scheme $S$ the
$X$-morphisms $S\to X'$ correspond bijectively to homomorphisms of
$\cO_X$-algebras $\mathcal A'\to \pi_*\cO_S$, where $\pi:S\to X$ is the
structure morphism; this is the defining adjunction of the relative
spectrum.  Applying $h'_*$ to the comorphism of $g'$ gives an
$\cO_X$-algebra map
\[
 \mathcal A'=h'_*\cO_{X'}\longrightarrow h'_*g'_*\cO_Y=f_*\cO_Y,
\]
hence, by the adjunction with $S=T$ and $\pi=h$, an $X$-morphism
$c:T\to X'$.  Composing with $g$ and using $h_*\cO_T=f_*\cO_Y$ shows
that $c\circ g$ and $g'$ correspond to the same algebra map
$\mathcal A'\to f_*\cO_Y$, so $c\circ g=g'$.  Uniqueness is the
injectivity half of the same adjunction: an $X$-morphism $c:T\to X'$
with $c\circ g=g'$ determines, and is determined by, the algebra map just
written, because $g^\#:\cO_T\to g_*\cO_Y$ is an isomorphism.

\cref{itm:ch42-su:unique}  If $g'_*\cO_Y=\cO_{X'}$, the displayed algebra
map is
$h'_*\cO_{X'}\to h'_*g'_*\cO_Y=h'_*\cO_{X'}$, the identity.  Thus $c$ is
the morphism induced by an isomorphism of $\cO_X$-algebras between
$\mathcal A'$ and $f_*\cO_Y$, hence is an isomorphism of relative
spectra.

\cref{itm:ch42-su:finite}  If $f$ is affine, properness makes it finite.  Apply
\cref{itm:ch42-su:unique} to the trivial factorization
$f=f\circ\mathrm{id}_Y$, with $X'=Y$, $g'=\mathrm{id}$; the hypothesis
$g'_*\cO_Y=\cO_Y$ is automatic and the conclusion is that
$c:T\to Y$ is an isomorphism inverse to $g$.  Conversely, if $g$ is an
isomorphism then $f=h\circ g$ is affine, being finite.
\end{proof}

\begin{remark}[What the universal property buys]
\label{rem:ch42-universal-buys}
Two consequences are used below without further comment.  The factorization
is unique up to unique isomorphism, so ``the finite part of $f$'' and ``the
birational part of $f$'' are well defined.  It is \emph{not} symmetric: the
universal property characterizes the largest
affine factor \emph{under} $Y$, and there is no corresponding statement
producing a smallest finite factor \emph{over} $X$.  This asymmetry is the
formal reason why the descent of \Cref{sec:ch42-descent} must run birational
first, and it is illustrated in \Cref{ex:ch42-second-factorization}.
\end{remark}

\subsection{The middle term is a normal integral projective variety}
\label{sec:ch42-middle}

\begin{lemma}\label{lem:bj-alteration-middle}
In \eqref{eq:bj-alteration-stein}, the scheme $T$ is a normal integral
projective variety, the map $h$ is finite surjective, and the map $g$ is
projective birational.
Moreover,
\[
                         K(T)=K(Y).
\tag{42.4}\label{eq:bj-alteration-function-fields}
\]
\end{lemma}

\begin{proof}
The source $Y$ is smooth, hence normal and integral.
The relative-normalization result
\cite[\href{https://stacks.math.columbia.edu/tag/035L}{Tag 035L}]
{stacks-project} says that the normalization of $X$ in a normal source is a
disjoint union of normal integral schemes and that its irreducible components
correspond to those of the source.
Since $Y$ has one irreducible component, $T$ is normal and integral.
The hypotheses of the tag are met: $f$ is proper, hence quasi-compact and
quasi-separated; $Y$ is normal; and $Y$ is noetherian, so every
quasi-compact open of $Y$ has finitely many irreducible components.  This
is exactly the checklist recorded in \Cref{app:stacks} for that tag, and it
is the reason normality of the middle term is not part of the Stein theorem
itself but of a second, independent statement.

The morphism $h$ is finite by Stein factorization.
A finite morphism to the projective $k$-variety $X$ is projective, and its
source $T$ is projective over $k$.
The image of $h$ is closed.
Since $f=h\circ g$ is surjective, the image is all of $X$, so $h$ is
finite surjective.

It remains to identify the generic fibre.
Let $\eta=\Spec K(X)$ be the generic point of $X$.
Generic finiteness of $f$ says that $Y_\eta$ is zero-dimensional, while
integrality says that its generic coordinate algebra is $K(Y)$; by
\Cref{lem:ch42-generic-finite}\,\cref{itm:ch42-gf:fibre}, in fact
$Y_\eta=\Spec K(Y)$.
After shrinking $X$ around $\eta$, the map is finite.
Consequently
\[
       (f_*\cO_Y)_\eta=H^0(Y_\eta,\cO_{Y_\eta})=K(Y).
\]
Taking the relative spectrum in \eqref{eq:bj-alteration-stein} gives
$T_\eta=\Spec K(Y)$, and hence $K(T)=K(Y)$.
Thus $g$ induces an isomorphism of function fields and is birational.

Finally $g$ is projective.
Choose a closed immersion $Y\hookrightarrow\mathbf P^N_X$ exhibiting the
projectivity of $f$.
The pair consisting of this immersion and $g$ gives a morphism
\[
              Y\longrightarrow\mathbf P^N_X\times_XT
                       =\mathbf P^N_T.
\]
Its graph is closed because $T\to X$ is separated, and the original map to
$\mathbf P^N_X$ is a closed immersion.
Hence $Y\to\mathbf P^N_T$ is a closed immersion, proving that $g$ is
projective.
\end{proof}

The identification of $T$ can be made completely explicit, and the explicit
form is what one computes with in examples.

\begin{lemma}
\label{lem:ch42-middle-normalization}
Let $f:Y\to X$ be a proper surjective generically finite morphism of
integral $k$-varieties with $Y$ normal, and let $T$ be as in
\eqref{eq:bj-alteration-stein}.  Then for every nonempty affine open
$U\subseteq X$,
\[
 \cO_T(h^{-1}U)
 =\{\,\xi\in K(Y):\xi\ \text{is integral over}\ \cO_X(U)\,\},
\]
the integral closure of $\cO_X(U)$ in $K(Y)$.  In particular $T$
depends only on $X$ and on the subfield $K(X)\subseteq K(Y)$, not on
$Y$; two alterations of $X$ with the same function field have the same
Stein middle term.
\end{lemma}

\begin{proof}
Write $A:=\cO_X(U)$ and $B:=\cO_T(h^{-1}U)=(f_*\cO_Y)(U)
=\cO_Y(f^{-1}U)$.  Since $h$ is finite, $B$ is a finite $A$-module,
hence integral over $A$; and $B\subseteq K(T)=K(Y)$ by
\eqref{eq:bj-alteration-function-fields}, because $T$ is integral.  Thus
$B$ is contained in the integral closure $\overline A$ of $A$ in
$K(Y)$.

Conversely let $\xi\in K(Y)$ be integral over $A$.  Then $\xi$ is
integral over $\cO_Y$ on $f^{-1}U$, and $Y$ is normal, so $\xi$ is a
regular function on $f^{-1}U$; that is, $\xi\in\cO_Y(f^{-1}U)=B$.  Hence
$\overline A\subseteq B$.

For the last sentence, note that the description just given involves only
$A$ and the field extension $K(X)\subseteq K(Y)$.
\end{proof}

\begin{lemma}
\label{lem:ch42-h-degree}
In \eqref{eq:bj-alteration-stein}, with $f$ an alteration of degree
$d=[K(Y):K(X)]$,
\[
 \deg g=1,\qquad \deg h=d,
\]
and $h$ is finite surjective of degree $d$.  Moreover:
\begin{enumerate}
\item\label{itm:ch42-hd:g-iso} $g$ is an isomorphism if and only if $f$
      is finite;
\item\label{itm:ch42-hd:h-iso} if $X$ is normal, then $h$ is an
      isomorphism if and only if $f$ is birational, that is, $d=1$.
\end{enumerate}
\end{lemma}

\begin{proof}
The degree statements are
\Cref{lem:ch42-degree-composite}\,\cref{itm:ch42-deg:comp} together with
\eqref{eq:bj-alteration-function-fields}: $K(T)=K(Y)$ gives
$\deg g=1$, and $d=\deg f=\deg g\cdot\deg h=\deg h$.

\cref{itm:ch42-hd:g-iso}  A finite morphism is affine, so this is
\Cref{lem:ch42-stein-universal}\,\cref{itm:ch42-su:finite} in one direction;
in the other, if $g$ is an isomorphism then $f=h\circ g$ is finite.

\cref{itm:ch42-hd:h-iso}  If $f$ is birational then $d=1$, so
$K(T)=K(X)$; by \Cref{lem:ch42-middle-normalization}, $\cO_T(h^{-1}U)$
is the integral closure of $\cO_X(U)$ in $K(X)$, which is
$\cO_X(U)$ because $X$ is normal.  Hence $h_*\cO_T=\cO_X$ and $h$,
being finite, is an isomorphism.  Conversely, if $h$ is an isomorphism then
$d=\deg h=1$.
\end{proof}

\begin{remark}[Connected fibres are not the birational argument]
\label{rem:ch42-connected-not-birational}
The connectedness of the fibres of $g$ is part of Stein factorization, but
it does not by itself imply birationality.
The decisive fact is \eqref{eq:bj-alteration-function-fields}.
The generic fibre has one point with residue field $K(Y)$, so the middle
variety and the source have the same function field.

The distinction is not pedantic.  A finite morphism of degree $d>1$ that is
purely inseparable with normal target --- as for the Frobenius
of a normal $X$ in \Cref{ex:ch42-frobenius-alteration}, or the map
$h$ of \Cref{ex:ch42-inseparable-cone} --- is a universal homeomorphism, so all its
fibres are single points and are certainly connected, and yet it is not birational when $d>1$.  Connectedness of
fibres is a topological statement; birationality is a statement about
function fields; and in positive characteristic the two come apart
completely.  What Stein factorization delivers is the algebraic statement
$g_*\cO_Y=\cO_T$, and it is that statement, through the computation of the
generic fibre, which yields \eqref{eq:bj-alteration-function-fields}.
\end{remark}

\begin{remark}[Why normality of the middle term is recorded]
\label{rem:ch42-normal-middle}
Both descent theorems used below require normal source and target.
Normality of $T$ is therefore not cosmetic.
It comes from the relative-normalization theorem and the normality of $Y$,
not merely from connectedness of the fibres.

It is worth seeing what would go wrong without it.
\Cref{thm:bj-birational-ep-descent} is stated for a projective birational
morphism $\beta:Y\to X$ of \emph{normal} integral projective varieties,
and normality is used on both sides: on the \emph{source} $Y$ through
\Cref{lem:bj-canonical-carrier}, which is applied on $Y$ with
$\eta=\beta^*\theta$, and on the \emph{target} $X$ through the standing
hypotheses of \Cref{thm:bj-carrier-reflection}.
\Cref{thm:bj-finite-normal-descent} is
stated for a finite surjective morphism $h:T\to X$ of normal integral
projective varieties, and there too both sides are used: normality of the
\emph{target} $X$ gives the identity $B^G=A$ of
\Cref{lem:bj-finite-normal-open}, which is exactly the statement that $X$
is integrally closed in $K(X)=M^G$, while normality of the \emph{source}
$T$ is what \Cref{lem:bj-canonical-carrier} needs when it is applied
upstairs.  \Cref{ex:ch41-node-not-open} is the reminder that the
target-side hypothesis is not decorative: there the source is the normal
curve $\mathbf P^1$ and the non-normal target is a nodal cubic, and the
analytification of the map fails to be open.
Thus $T$ has to be normal for two different reasons in two different
chapters, and neither reason is supplied by
\Cref{thm:appB-stein-factorization} alone.
\end{remark}

\subsection{The separation, proved}
\label{sec:ch42-separation}

We can now state the separation of the two mechanisms as a single
proposition.  Everything in it has been proved above; the point of
collecting it is that the hypotheses of
\Cref{thm:bj-birational-ep-descent} and of
\Cref{thm:bj-finite-normal-descent} are matched item by item, so that the
descent in \Cref{sec:ch42-descent} can proceed by quotation.

\begin{proposition}
\label{prop:ch42-separation}
Let $X$ be a normal integral projective $k$-variety and let
$f:Y\to X$ be a proper surjective generically finite morphism of degree
$d$ with $Y$ normal, integral and projective.  Let
$Y\xrightarrow{g}T\xrightarrow{h}X$ be the Stein factorization
\eqref{eq:bj-alteration-stein}.  Then:
\begin{enumerate}
\item\label{itm:ch42-sep:T} $T$ is a normal integral projective
      $k$-variety with $K(T)=K(Y)$;
\item\label{itm:ch42-sep:g} $g:Y\to T$ is a projective birational morphism
      of normal integral projective $k$-varieties, so that
      \Cref{thm:bj-birational-ep-descent} applies to $g$ and to every
      class in $N^1(T)_{\mathbf R}$;
\item\label{itm:ch42-sep:h} $h:T\to X$ is a finite surjective morphism of
      degree $d$ of normal integral projective $k$-varieties, so that
      \Cref{thm:bj-finite-normal-descent} applies to $h$ and to every
      class in $N^1(X)_{\mathbf R}$;
\item\label{itm:ch42-sep:canonical} the factorization is unique up to unique
      isomorphism among factorizations $f=h'\circ g'$ with $h'$ finite
      and $g'_*\cO_Y=\cO_{X'}$;
\item\label{itm:ch42-sep:degenerate} $g$ is an isomorphism precisely when
      $f$ is finite, and $h$ is an isomorphism precisely when $f$ is
      birational.
\end{enumerate}
\end{proposition}

\begin{proof}
\cref{itm:ch42-sep:T} and the birationality and projectivity in
\cref{itm:ch42-sep:g} are \Cref{lem:bj-alteration-middle}; the proof given
there used smoothness of $Y$ only through its normality and integrality,
which are hypotheses here.  Normality of $Y$ is a hypothesis and normality
of $T$ is part of \cref{itm:ch42-sep:T}, so the source and target of $g$
are as \Cref{thm:bj-birational-ep-descent} requires; that theorem is stated
for an arbitrary real class on the target, with no positivity hypothesis, by
\Cref{rem:ch39-downstream}.

\cref{itm:ch42-sep:h}  Finiteness and surjectivity are
\Cref{lem:bj-alteration-middle}, the degree is
\Cref{lem:ch42-h-degree}, and normality, integrality and projectivity of
both $T$ and $X$ are hypotheses or have just been proved.  These are
exactly the standing hypotheses of \Cref{thm:bj-finite-normal-descent},
which likewise imposes no positivity on the class, by
\Cref{rem:ch41-descent-hypothesis-map}.

\cref{itm:ch42-sep:canonical} is
\Cref{lem:ch42-stein-universal}\,\cref{itm:ch42-su:unique}, and
\cref{itm:ch42-sep:degenerate} is \Cref{lem:ch42-h-degree}.
\end{proof}

\begin{warning}{The Stein factorization is not the flattening square}
\label{warn:ch42-stein-not-flattening}
Two commutative diagrams built out of a finite or generically finite
morphism appear in this part of the book, and they are easy to conflate
because both involve a birational change and a finite morphism.  They are
different.

The \emph{Stein factorization} \eqref{eq:bj-alteration-stein} of an
alteration $f:Y\to X$ is a factorization of $f$: a triangle, not a
square.  Its birational factor $g$ sits \emph{above} its finite factor
$h$ and contracts part of the source $Y$.
Stein factorization exists for every proper morphism in the Noetherian
setting; the first factor is birational here because $f$ is generically
finite.
The \emph{flattening square} \eqref{eq:bj-flatification-square} of
\Cref{lem:bj-finite-flatification} is applied to a finite morphism
$h:T\to X$ that is already given; it produces birational modifications
$\beta:Z\to X$ of the target \emph{and} $\alpha:W\to T$ of the source,
together with a finite \emph{flat} $q:W\to Z$, and it exists by
Raynaud--Gruson flattening.  Its purpose is to make a determinant available,
not to separate mechanisms.

In the proof of \Cref{thm:bj-normal-envelope} both occur, one inside the
other: the Stein factorization is formed first, and the flattening square is
formed afterwards, inside \Cref{thm:bj-finite-normal-descent}, for the
finite factor $h$.  \Cref{ex:ch42-cone-stein} and
\Cref{ex:ch41-flatten-cone} compute the two diagrams for the same surface,
and they are visibly different: the Stein middle term is $\mathbf P^2$,
while the flattening replaces $\mathbf P^2$ by $\mathbf F_1$ and
$\mathbf P(1,1,2)$ by $\mathbf F_2$.
\end{warning}

\section{The class ledger}
\label{sec:ch42-ledger}

The descent moves a psh function down two morphisms, and at each stage the
function must be psh for a definite numerical class.  Keeping track of those
classes is elementary, but there is one arithmetical trap: the finite step
of \Cref{chap:bj-finite-descent} contains a division by the degree, and it
is tempting to imagine that the degree has to be inserted somewhere else to
compensate.  It does not.  This section records the exact bookkeeping.

\subsection{Functoriality and the two projection formulae}
\label{sec:ch42-functoriality}

Fix $\theta\in N^1(X)_{\mathbf R}$.
Functoriality of numerical pullback in the factorization $f=h\circ g$
gives
\[
                         f^*\theta=g^*h^*\theta.
\tag{42.5}\label{eq:bj-alteration-class-ledger}
\]
This elementary identity controls the entire proof.
The smooth theorem applies on $Y$ in the class on the left.
Birational descent along $g$ returns the class $h^*\theta$ on $T$.
Finite descent along $h$ then returns $\theta$ on $X$.

It would be incorrect to divide $f^*\theta$ by the generic degree at this
stage.
The degree normalization belongs to the determinant-norm construction inside
finite descent.
The birational step carries no degree factor.

\begin{lemma}
\label{lem:ch42-class-ledger}
Let $f:Y\to X$ be a proper surjective generically finite morphism of
degree $d$ between integral projective $k$-varieties of dimension $n$,
with Stein factorization $f=h\circ g$ as in
\eqref{eq:bj-alteration-stein}.  Then:
\begin{enumerate}
\item\label{itm:ch42-cl:functor}
      $f^*=g^*\circ h^*$ as maps
      $N^1(X)_{\mathbf R}\to N^1(Y)_{\mathbf R}$, and each of $f^*$,
      $g^*$, $h^*$ is $\mathbf R$-linear and continuous for the
      finite-dimensional topologies;
\item\label{itm:ch42-cl:degree}
      for all $\theta_1,\ldots,\theta_n\in N^1(X)_{\mathbf R}$,
      \[
       (f^*\theta_1\cdots f^*\theta_n)=d\,(\theta_1\cdots\theta_n),
       \qquad
       (g^*\eta_1\cdots g^*\eta_n)=(\eta_1\cdots\eta_n)
      \]
      for $\eta_j\in N^1(T)_{\mathbf R}$, and
      $(h^*\theta_1\cdots h^*\theta_n)=d\,(\theta_1\cdots\theta_n)$;
\item\label{itm:ch42-cl:inject}
      $f^*$, $g^*$ and $h^*$ are injective;
\item\label{itm:ch42-cl:nef}
      for $\theta\in N^1(X)_{\mathbf R}$, $\theta$ is nef if and
      only if $f^*\theta$ is nef; the analogous assertion holds for
      $h$ and for every $\eta\in N^1(T)_{\mathbf R}$ along $g$;
\item\label{itm:ch42-cl:psef}
      if $\theta$ is pseudoeffective then so is $f^*\theta$.
\end{enumerate}
\end{lemma}

\begin{proof}
\cref{itm:ch42-cl:functor}  Pullback of Cartier divisors is functorial for
composites of morphisms of schemes.  To check numerical equivalence,
apply the cycle projection formula \eqref{eq:appF-projection} to an
integral curve $C$ in the source: a contracted curve has pullback degree
zero, while a curve mapping onto $D$ has degree
$[K(C):K(D)]\deg_D(L)$ for a line bundle $L$ on the target.
A numerically trivial $L$ therefore remains numerically trivial.
Linearity and continuity follow on finite-dimensional real vector spaces.

\cref{itm:ch42-cl:degree} is \Cref{cor:appF-generically-finite}.
For the stated integral scope, $T$ is integral because its coordinate
ring over a nonempty affine open of $X$ is the domain of regular
functions on the corresponding open of $Y$.
The generic-fibre calculation in \Cref{lem:bj-alteration-middle}
does not require normality: it gives $K(T)=K(Y)$, so $g$ has
degree one and $h$ has degree $d$.  Apply the intersection formula
to each of $f,g,h$.  These are cycle intersection statements;
no arbitrary pushforward of Cartier classes on $N^1$ is asserted.

\cref{itm:ch42-cl:inject}  It suffices to treat $f$, the other two being
special cases.  Let $\theta\in N^1(X)_{\mathbf R}$ with $f^*\theta=0$,
and let $C\subset X$ be an integral curve.  Since $f$ is proper and
surjective, $f^{-1}(C)$ is closed and surjects onto $C$; choose an
irreducible component $V$ of $f^{-1}(C)$ with $f(V)=C$, and let
$C'\subseteq V$ be an integral curve dominating $C$.
Intersect $V$ with $\dim V-1$ general members of a very ample
linear system.  Over the generic point of $C$ this cuts a nonempty
zero-dimensional scheme, by positivity of ample top degree on the
projective generic fibre.  Choose a component of the resulting
one-cycle that dominates $C$.  The projection formula \eqref{eq:appF-projection} gives
\[
 0=(f^*\theta\cdot C')=(\theta\cdot f_*C')=e\,(\theta\cdot C),
\]
where $e=[K(C'):K(C)]\ge1$.  Hence $(\theta\cdot C)=0$ for every
integral curve $C$, that is, $\theta=0$ in $N^1(X)_{\mathbf R}$ by
\Cref{def:appF-numerical}.

\cref{itm:ch42-cl:nef}  If $\theta$ is nef and $C'\subset Y$ is an
integral curve, then either $f(C')$ is a point, in which case
$(f^*\theta\cdot C')=0$, or $f(C')=C$ is an integral curve and
$(f^*\theta\cdot C')=e(\theta\cdot C)\ge0$ as above.  Conversely, if
$f^*\theta$ is nef, the same computation with $C'$ chosen over a given
$C$ gives $e(\theta\cdot C)\ge0$, hence $(\theta\cdot C)\ge0$.

\cref{itm:ch42-cl:psef}  Use the closed cone generated by effective rational Cartier classes.
After clearing denominators, pullback of an effective Cartier divisor
along the dominant morphism is effective: its local equation pulls
back to a nonzero regular function on the integral source.
Taking nonnegative real combinations and then limits proves the
assertion, by continuity of $f^*$.
\end{proof}

\begin{warning}{The birational step carries no degree factor}
\label{warn:ch42-no-degree}
The degree normalizations used in the proof of
\Cref{thm:bj-normal-envelope} occur inside
\Cref{thm:bj-finite-normal-descent}: once as the real normalizing factor
$1/d$ in the fibre average \eqref{eq:ch41-average-shape}, and once as the
rational factor $1/d$ in the class identity \eqref{eq:bj-norm-class}.  It
does not occur in \Cref{thm:bj-birational-ep-descent}, and it must not be
inserted there.  A reader who writes $f^*\theta/d$ for the class upstairs
--- an instinct imported from the theory of finite covers of curves, where
degrees of line bundles do get divided --- will find that the classes fail
to match at the middle term, because $g^*$ preserves intersection numbers
while $h^*$ multiplies them by $d$, by
\Cref{lem:ch42-class-ledger}\,\cref{itm:ch42-cl:degree}.

The point is also visible in \Cref{rem:ch41-division-by-d}: the two
divisions by $d$ are a division of a real number and a division of a
numerical class, neither of them a division inside $k$ or inside a finite
algebra.  They therefore survive $\operatorname{char}k\mid d$, which is
exactly the situation of \Cref{ex:ch42-inseparable-cone}.
\end{warning}

\subsection{The descent chain}
\label{sec:ch42-descent-chain}

The dependency can be displayed as
\[
\begin{tikzcd}[column sep=large]
 \operatorname{EP}_Y(g^*h^*\theta)
   \arrow[r,"\text{birational descent along }g"]
 &\operatorname{EP}_T(h^*\theta)
   \arrow[r,"\text{finite descent along }h"]
 &\operatorname{EP}_X(\theta).
\end{tikzcd}
\tag{42.6}\label{eq:bj-alteration-descent-chain}
\]
Every arrow in \eqref{eq:bj-alteration-descent-chain} points from a theorem
already established to a new conclusion.
There is no appeal to the normal case in the proof of either arrow.
This is the circularity firewall for the present chapter.

Three features of the chain deserve to be stated explicitly, because each
is a hypothesis that some plausible variant of the argument would violate.

\emph{The middle class is not assumed positive.}  The class $h^*\theta$
appearing at the middle term is an arbitrary real class on $T$; it is nef
exactly when $\theta$ is, by
\Cref{lem:ch42-class-ledger}\,\cref{itm:ch42-cl:nef}, and $\theta$ is not
assumed nef.  This is why \Cref{chap:bj-reflection} was proved without a
positivity hypothesis, as \Cref{rem:ch39-downstream} records; a version of
birational descent stated only for nef or ample classes would be useless
here.

\emph{The arrows are implications, not equivalences.}  Both
\Cref{thm:bj-birational-ep-descent} and
\Cref{thm:bj-finite-normal-descent} are one-way as proved, by
\Cref{rem:ch39-one-way} and \Cref{warn:ch41-one-way}.  The chain therefore
has to be traversed from left to right, and the leftmost hypothesis has to
be supplied from outside.  It is supplied by
\Cref{thm:bj-smooth-envelope}, which is the reason the alteration must have
a \emph{smooth} source and not merely a normal one.

\emph{The chain is finite.}  There is exactly one birational step and
exactly one finite step, and no limit or iteration is involved.  This is
worth saying because the analogous reduction in some treatments of
resolution proceeds by induction on a resolution invariant; nothing of the
kind happens here, and in particular no compatibility between different
alterations of $X$ is needed.  Independence of the choice of alteration is
a \emph{consequence} of \Cref{thm:bj-normal-envelope}, recorded in
\Cref{cor:ch42-independence-of-alteration}, not an input to it.

\begin{remark}[Real classes and the drift]
\label{rem:ch42-real-classes}
The class $\theta$ is real and is never approximated by rational classes
in this chapter.  This is deliberate.  The reduction from real to rational
classes, and the $O(m^{-1})$ drift that it produces, are carried out once
and for all inside \Cref{chap:bj-smooth-characteristic}
(\Cref{sec:ch40-reduction} and \Cref{lem:ch40-drift-nef-monotone}); by the
time \Cref{thm:bj-smooth-envelope} is quoted, the real case is already
proved.  Similarly, the varying classes of \Cref{def:bj-psh} are handled
inside \Cref{chap:bj-finite-descent}, through
\Cref{cor:bj-coefficient-norm-limit} and
\Cref{cor:ch41-real-weights}.  The present chapter therefore performs no
approximation of classes at all, and \eqref{eq:bj-alteration-class-ledger} is
an exact identity, not an identity up to $O(m^{-1})$.  This is the reason
the chapter can be short in its logical core while the two chapters it
quotes cannot.
\end{remark}

\section{Descent from a smooth alteration}
\label{sec:ch42-descent}

We now prove the implication that the chapter exists to prove.  It is
stated for a normal source, because that is all the argument uses; the
smoothness of de Jong's alteration is needed only to supply the hypothesis
$\operatorname{EP}_Y$, through \Cref{thm:bj-smooth-envelope}, and it plays
no other role.  Separating the two makes the logical structure visible and
yields the tower statements of \Cref{sec:ch42-descent-consequences} for
free.

\subsection{The descent theorem}
\label{sec:ch42-descent-theorem}

\begin{theorem}
\label{thm:ch42-alteration-descent}
Let $X$ and $Y$ be normal integral projective $k$-varieties, let
\[
 f:Y\longrightarrow X
\]
be a proper surjective generically finite morphism, and let
$\theta\in N^1(X)_{\mathbf R}$.  Then
\[
 \operatorname{EP}_Y(f^*\theta)
 \quad\Longrightarrow\quad
 \operatorname{EP}_X(\theta).
\]
No hypothesis is imposed on the degree of $f$, on the separability of
$K(Y)/K(X)$, on the characteristic of $k$, on the dimension, or on the
positivity of $\theta$.
\end{theorem}

\begin{proof}
The proof has three steps: factor, descend birationally, descend finitely.
Each step is an application of a theorem proved in an earlier chapter, and
the work consists in verifying that its hypotheses hold.

\medskip
\noindent\emph{Step 1: factor.}
Let
\[
 Y\xrightarrow{\ g\ }T\xrightarrow{\ h\ }X
\]
be the Stein factorization \eqref{eq:bj-alteration-stein} of $f$, and put
$d:=\deg f$.  By \Cref{prop:ch42-separation}:
$T$ is a normal integral projective $k$-variety with $K(T)=K(Y)$;
$g$ is a projective birational morphism of normal integral projective
$k$-varieties; and $h$ is a finite surjective morphism of degree $d$
of normal integral projective $k$-varieties.  All the hypotheses of that
proposition are satisfied: $X$ and $Y$ are normal integral projective by
assumption, and $f$ is proper, surjective and generically finite.

Set
\[
 \eta:=h^*\theta\in N^1(T)_{\mathbf R}.
\]
By \Cref{lem:ch42-class-ledger}\,\cref{itm:ch42-cl:functor},
\[
 g^*\eta=g^*h^*\theta=f^*\theta,
\]
which is \eqref{eq:bj-alteration-class-ledger}.  The hypothesis of the
theorem is therefore exactly $\operatorname{EP}_Y(g^*\eta)$.

\medskip
\noindent\emph{Step 2: descend birationally, from $Y$ to $T$.}
Apply \Cref{thm:bj-birational-ep-descent} to the projective birational
morphism $g:Y\to T$ of normal integral projective $k$-varieties and to
the class $\eta\in N^1(T)_{\mathbf R}$.  Its hypotheses are met: the
source and target are normal integral projective by Step~1, the morphism is
projective birational by Step~1, the class is an arbitrary element of
$N^1(T)_{\mathbf R}$, which is what that theorem allows --- it assumes no
positivity, precisely in order to be applicable here, as
\Cref{rem:ch39-downstream} records --- and
$\operatorname{EP}_Y(g^*\eta)$ is the hypothesis of the present theorem.
The conclusion is
\[
 \operatorname{EP}_T(\eta)=\operatorname{EP}_T(h^*\theta).
\]

We pause on one point that the notation hides.  The envelope property
$\operatorname{EP}_T(\eta)$ is a statement about \emph{all} bounded-above
nonempty families in $\PSH_T(\eta)$, not only about families pulled back
from $X$.  This is essential because finite descent requires the full
envelope property on its source, not a restricted assertion about a
preselected collection of families.  Birational descent obtains that full
statement by pulling an arbitrary $T$-family to $Y$.  Both descent theorems
are stated in the strong form ``every family'' for exactly this reason.

\medskip
\noindent\emph{Step 3: descend finitely, from $T$ to $X$.}
Apply \Cref{thm:bj-finite-normal-descent} to the finite surjective morphism
$h:T\to X$ of normal integral projective $k$-varieties and to the class
$\theta\in N^1(X)_{\mathbf R}$.  Its hypotheses are met: source and target
are normal integral projective by Step~1, $h$ is finite surjective by
Step~1, and the class is arbitrary --- that theorem too assumes no
positivity, by \Cref{rem:ch41-descent-hypothesis-map}.  Its hypothesis
$\operatorname{EP}_T(h^*\theta)$ is the conclusion of Step~2.  The
conclusion is $\operatorname{EP}_X(\theta)$, which is what was to be
proved.

\medskip
\noindent\emph{No circularity.}  The initial envelope hypothesis is on $Y$.
Step~2 consumes it to prove EP on $T$, and Step~3 consumes that derived
EP to prove EP on $X$.  At no point was $\operatorname{EP}$ asserted for a variety
for which it had not been either assumed or previously derived; in
particular \Cref{thm:bj-final-envelope} was not used, in accordance with
\Cref{warn:ch42-firewall}.
\end{proof}

\begin{remark}[Where each hypothesis of \Cref{thm:ch42-alteration-descent}
is used]
\label{rem:ch42-descent-hypothesis-map}
Normality of $Y$ is used once, in \Cref{prop:appB-stein-normal} through
\Cref{lem:bj-alteration-middle}, to make the middle term $T$ normal; it is
used a second time, implicitly, because
\Cref{thm:bj-birational-ep-descent} requires a normal source.  Normality of
$X$ is required by finite descent in Step~3, through
\Cref{lem:bj-finite-normal-open} and through the standing hypotheses of
\Cref{thm:bj-carrier-reflection} inside the proof of
\Cref{thm:bj-finite-normal-descent}; see
\Cref{rem:ch41-descent-hypothesis-map}.  Integrality of both is used for the
existence of a single function field, hence for the notion of degree, and
for generic finiteness of Fubini--Study functions.  Projectivity is used for
compactness of the analytifications, for the existence of centres
(\Cref{prop:centers-proper-separated}), for the projectivity of $g$ in
\Cref{lem:bj-alteration-middle}, and for the flattening construction inside
\Cref{chap:bj-finite-descent}.  Properness of $f$ is used for the Stein
factorization, surjectivity for the surjectivity of $h$, and generic
finiteness for \eqref{eq:bj-alteration-function-fields}.

In particular the theorem imposes no additional condition on, and
needs nothing from, the exceptional locus of $f$, the ramification divisor
of $h$, or the singularities of $X$ beyond normality.
\end{remark}

\begin{warning}{The order supplied by Stein factorization}
\label{warn:ch42-order}
The proved argument runs birational descent first and finite descent
second because Stein factorization places the maps in that order.
It does not assert that the reverse order is logically invalid.

A reverse route needs a factorization
$Y\xrightarrow{q}Z\xrightarrow{\beta}X$ with $q$ finite and
$\beta$ projective birational.  Such a factorization of the given
source is not furnished by Stein factorization.
When it exists and $Z$ is normal, both descent theorems apply
in that order, as \Cref{ex:ch42-second-factorization} shows.

A nonnormal middle term can also be repaired.  Since $Y$ is normal,
the finite dominant map $Y\to Z$ lifts to the finite normalization
$Z^\nu\to Z$.  On affine charts the integral closure of $\cO_Z$
in $K(Z)$ maps into the normal finite source algebra, which remains
finite over that intermediate integral closure.  Thus $Y\to Z^\nu$
is finite surjective and $Z^\nu\to X$ is projective birational.
This uses ordinary finite normalization and the already available
normal-target descent theorems, without later unibranch descent.
Stein factorization supplies the required normal middle term
canonically for every alteration here.
\end{warning}

\begin{remark}[Why one cannot descend along $f$ in a single step]
\label{rem:ch42-not-one-step}
It is worth recording precisely which hypothesis of each descent theorem the
alteration $f$ itself fails, since ``$f$ is neither birational nor
finite'' is a statement about $f$ and not yet a statement about the
proofs.

\Cref{thm:bj-birational-ep-descent} is proved by reflecting a psh function
through \Cref{thm:bj-carrier-reflection}, whose entire mechanism --- Cartier
b-divisors, homogeneous Siu division, and the identification of divisorial
loci --- presupposes $K(Y)=K(X)$.  When $d>1$ there is no map
$X^{\divis}\to Y^{\divis}$ inverse to $\beta^{\An}$, and the common-field identification in
\Cref{lem:ch39-beta-an}\,\cref{itm:ch39-ban:val} cannot be invoked.

\Cref{thm:bj-finite-normal-descent} is proved by averaging over fibres with
the weights of \eqref{eq:ch41-weights}, and the very first requirement is
that a fibre be a finite scheme, so that
$\sum_{w\mid z}e(w/z)=d$ is a finite sum of positive integers.  An
alteration that contracts a curve has fibres that are not finite over the
base, and \Cref{ex:ch42-no-separation} exhibits a psh function whose values
on such a fibre are not constant, so that no finite average of them is
defined.

There is a third temptation, which is more interesting because it almost
works: the sheaf $f_*\cO_Y$ \emph{is} a finite $\cO_X$-algebra of
generic rank $d$, even though $f$ is not finite.  That does not
make it locally free; flattening is still needed for the constant-rank
determinant package.  But
$f_*\cO_Y=h_*g_*\cO_Y=h_*\cO_T$ by the definition of the Stein
factorization, so the algebra one is using is the algebra of $h$, and the
resulting averaging operator takes as its input a Fubini--Study function on
$T$, not on $Y$.  Getting from $Y$ to $T$ is exactly the birational
step.  The apparently one-step argument therefore contains the two-step
argument, with its birational half concealed inside the identification
$f_*\cO_Y=h_*\cO_T$.  This is worked out in
\Cref{ex:ch42-no-separation}\,(3).
\end{remark}

\subsection{Consequences of the descent theorem}
\label{sec:ch42-descent-consequences}

\begin{corollary}
\label{cor:ch42-alteration-tower}
Let $Z\xrightarrow{\ b\ }Y\xrightarrow{\ a\ }X$ be proper surjective
generically finite morphisms of normal integral projective $k$-varieties,
and let $\theta\in N^1(X)_{\mathbf R}$.  Then
\[
 \operatorname{EP}_Z\bigl((a\circ b)^*\theta\bigr)
 \ \Longrightarrow\
 \operatorname{EP}_Y(a^*\theta)
 \ \Longrightarrow\
 \operatorname{EP}_X(\theta).
\]
In particular, if $\operatorname{EP}_Z((a\circ b)^*\theta)$ holds for some
tower, it holds for every intermediate normal projective variety.
\end{corollary}

\begin{proof}
Apply \Cref{thm:ch42-alteration-descent} to $b$ with the class
$a^*\theta$, and then to $a$ with the class $\theta$; the classes
match because $(a\circ b)^*=b^*a^*$.  Both $a$ and $b$ are proper
surjective generically finite by
\Cref{lem:ch42-degree-composite}, and their sources and targets are normal
integral projective by hypothesis.
\end{proof}

\begin{corollary}
\label{cor:ch42-one-alteration}
Let $X$ be a normal integral projective $k$-variety and
$\theta\in N^1(X)_{\mathbf R}$.  Suppose that there exists \emph{some}
proper surjective generically finite morphism $f:Y\to X$ with $Y$ normal
integral projective such that $\operatorname{EP}_Y(f^*\theta)$ holds.
Then $\operatorname{EP}_X(\theta)$ holds.
\end{corollary}

\begin{proof}
This is \Cref{thm:ch42-alteration-descent} verbatim; it is stated separately
because it is the form in which the theorem is used, and because it makes
explicit that no compatibility between different alterations is required.
\end{proof}

\begin{corollary}
\label{cor:ch42-char-zero-shortcut}
Suppose $\operatorname{char}k=0$, let $X$ be a normal integral projective
$k$-variety and let $\theta\in N^1(X)_{\mathbf R}$.  Then
$\operatorname{EP}_X(\theta)$ holds, and the proof uses only
\Cref{thm:appB-resolution}, \Cref{thm:bj-smooth-envelope} and
\Cref{thm:bj-birational-ep-descent}; neither Stein factorization nor
\Cref{chap:bj-finite-descent} is needed.
\end{corollary}

\begin{proof}
By \Cref{thm:appB-resolution} there is a projective birational morphism
$\pi:\widetilde X\to X$ with $\widetilde X$ smooth and projective; it is
integral because $X$ is and $\pi$ is birational.  By
\Cref{thm:bj-smooth-envelope}, $\operatorname{EP}_{\widetilde X}
(\pi^*\theta)$ holds.  Since $\widetilde X$ and $X$ are normal integral
projective and $\pi$ is projective birational,
\Cref{thm:bj-birational-ep-descent} gives $\operatorname{EP}_X(\theta)$.
\end{proof}

\begin{remark}[What the shortcut costs and what it saves]
\label{rem:ch42-shortcut-scope}
\Cref{cor:ch42-char-zero-shortcut} is a genuine simplification, and a reader
whose interest is confined to characteristic zero may take it as the normal
case and pass directly to \Cref{chap:bj-subintegral-normalization}.  It
should be noted, though, what it does \emph{not} simplify.  It removes
\Cref{chap:bj-finite-descent} from the proof of the normal case only; the
finite descent theorem is still needed in
\Cref{chap:bj-envelope-consequences} whenever a genuinely finite morphism
occurs, and the determinant-norm machinery is used again in
\Cref{cor:ch42-finite-bijection}.  It also does not remove
\Cref{chap:bj-smooth-characteristic}: the characteristic-zero half of that
chapter, with multiplier ideals in place of test ideals, is exactly what
supplies $\operatorname{EP}_{\widetilde X}$.
\end{remark}

\section{The envelope property on normal varieties}
\label{sec:ch42-normal-theorem}

\subsection{The theorem}
\label{sec:ch42-theorem}

\begin{theorem}\label{thm:bj-normal-envelope}
Let $X$ be a normal integral projective variety over an algebraically
closed trivially valued field $k$.
Then
\[
                         \operatorname{EP}_X(\theta)
\tag{42.7}\label{eq:bj-normal-envelope}
\]
holds for every $\theta\in N^1(X)_{\mathbf R}$, in arbitrary dimension and
characteristic.
\end{theorem}

\begin{proof}
Apply \Cref{thm:bj-de-jong-input} with $Z=\varnothing$.  Because $k$ is
perfect, its regular source is smooth; because $k$ is algebraically closed,
the finite ground-field extension in de Jong's construction is trivial.
The properness argument in the proof of that theorem therefore supplies a
smooth integral projective alteration $f:Y\to X$.

Fix $\theta\in N^1(X)_{\mathbf R}$.  The smooth theorem gives
\[
                         \operatorname{EP}_Y(f^*\theta).
\tag{42.8}\label{eq:bj-alteration-smooth-ep}
\]
Now apply \Cref{thm:ch42-alteration-descent} to $f$.  In its Stein
factorization $f=h\circ g$, its birational step gives
\[
                         \operatorname{EP}_T(h^*\theta).
\tag{42.9}\label{eq:bj-alteration-middle-ep}
\]
The finite step then gives $\operatorname{EP}_X(\theta)$.  The class is
arbitrary: the smooth theorem and both descent theorems are stated for real
classes without a positivity assumption.  Finally, the three inputs just
cited are independent of \Cref{thm:bj-final-envelope}, so the firewall of
\Cref{warn:ch42-firewall} is respected.
\end{proof}

\begin{remark}[The three quotations, and what each contributes]
\label{rem:ch42-three-quotations}
It is worth separating the three imported statements by what they are
responsible for.

\Cref{thm:bj-de-jong-input} is responsible for \emph{existence}.
De Jong's theorem is an imported foundational input.
Stein factorization, relative normalization, flattening and the earlier
ideal/Siu packages also retain their stated external inputs;
this chapter verifies the assembly hypotheses.

\Cref{thm:bj-smooth-envelope} supplies the initial envelope property
on the smooth source, using asymptotic multiplier or test ideals.
The subsequent descent invokes the already established carrier and norm
constructions at their stated scopes.

\Cref{thm:ch42-alteration-descent} is responsible for \emph{transport}, and
it is transport along a morphism that mixes two mechanisms.  The mixing is
undone by \Cref{thm:appB-stein-factorization}, which is why the geometric
input of this chapter is really two theorems and not one.
\end{remark}

\subsection{First consequences}
\label{sec:ch42-consequences}

\begin{corollary}\label{cor:bj-normalization-source-ep}
Let $X$ be any integral projective variety and let
$\nu:X^\nu\to X$ be its normalization.
For every $\theta\in N^1(X)_{\mathbf R}$,
\[
                  \operatorname{EP}_{X^\nu}(\nu^*\theta)
\tag{42.10}\label{eq:bj-normalization-source-ep}
\]
holds.
\end{corollary}

\begin{proof}
The normalization of a finite-type integral $k$-scheme is finite by
\cite[\href{https://stacks.math.columbia.edu/tag/035S}{Tag 035S}]
{stacks-project}, because such a scheme is Nagata.
Thus $X^\nu$ is integral, normal, and projective.
Apply \Cref{thm:bj-normal-envelope} to $X^\nu$ and the real class
$\nu^*\theta$.

In detail: $X^\nu$ is integral because $X$ is and normalization does not
change the function field; it is normal by construction; and it is
projective because $\nu$ is finite, hence projective, and $X$ is
projective, so that $X^\nu$ is projective over $k$.  The class
$\nu^*\theta$ lies in $N^1(X^\nu)_{\mathbf R}$ because $\nu$ is a
morphism of schemes and numerical pullback is defined for
$\mathbf R$-Cartier classes, by \Cref{prop:appF-mixed-cartier}.  No
hypothesis on $X$ beyond integrality and projectivity is used, and in
particular $X$ is \emph{not} assumed unibranch here; unibranchness enters
only when one wishes to descend \emph{from} $X^\nu$ \emph{to} $X$,
which is the subject of
\Cref{chap:bj-universal-homeomorphisms,chap:bj-subintegral-normalization}.
\end{proof}

\begin{corollary}
\label{cor:ch42-independence-of-alteration}
Let $X$ be a normal integral projective $k$-variety,
$\theta\in N^1(X)_{\mathbf R}$, and let $f_1:Y_1\to X$ and
$f_2:Y_2\to X$ be any two proper surjective generically finite morphisms
with $Y_1,Y_2$ normal integral projective.  Then
$\operatorname{EP}_{Y_1}(f_1^*\theta)$,
$\operatorname{EP}_{Y_2}(f_2^*\theta)$ and
$\operatorname{EP}_X(\theta)$ all hold, and in particular the conclusion
of \Cref{thm:ch42-alteration-descent} does not depend on the alteration
chosen.
\end{corollary}

\begin{proof}
$\operatorname{EP}_X(\theta)$ is \Cref{thm:bj-normal-envelope}.  Each
$Y_j$ is a normal integral projective $k$-variety, so
\Cref{thm:bj-normal-envelope} applies to it and to the class $f_j^*\theta$
as well.
\end{proof}

\begin{corollary}
\label{cor:ch42-quotient}
Let $Y$ be a normal integral projective $k$-variety, let $G$ be a
finite group acting on $Y$ by $k$-automorphisms with
$Y/G$ projective, and let $\theta\in N^1(Y/G)_{\mathbf R}$.  Then
$\operatorname{EP}_{Y/G}(\theta)$ holds.  No hypothesis relating
$|G|$ to $\operatorname{char}k$ is needed.
\end{corollary}

\begin{proof}
The quotient $X:=Y/G$ is a normal integral projective $k$-variety: the
quotient map $\pi:Y\to X$ is finite surjective, and $\cO_X=(\pi_*\cO_Y)^G$
is a ring of invariants of a normal domain, hence a normal domain.  Apply
\Cref{thm:bj-normal-envelope} to $X$.  Alternatively, apply
\Cref{thm:ch42-alteration-descent} to $\pi$, whose source has the envelope
property by \Cref{thm:bj-normal-envelope}.

The absence of a hypothesis on $|G|$ is the point.  A proof that averaged
over $G$ would need $|G|$ invertible in $k$; the determinant norm of
\Cref{sec:ch41-determinant} needs nothing of the kind, as
\Cref{rem:ch41-division-by-d} explains.
\end{proof}

\begin{corollary}
\label{cor:ch42-birational-bijection}
Let $\beta:Y\to X$ be a projective birational morphism of normal integral
projective $k$-varieties and let $\theta\in N^1(X)_{\mathbf R}$.  Then
\[
 \beta^*:\PSH_X(\theta)\longrightarrow\PSH_Y(\beta^*\theta)
\]
is a bijection.
\end{corollary}

\begin{proof}
By \Cref{thm:bj-normal-envelope} applied to the normal integral projective
variety $Y$ and the class $\beta^*\theta$, the hypothesis
$\operatorname{EP}_Y(\beta^*\theta)$ of
\Cref{thm:bj-birational-ep-descent} holds.  That theorem then gives
\eqref{eq:bj-ordinary-reflection-under-ep}, which is the assertion.  The
case $\PSH_Y(\beta^*\theta)=\varnothing$ is included, both sides being
empty by \Cref{thm:ch39-pullback}\,\cref{itm:ch39-pb:psh}.
\end{proof}

\begin{corollary}
\label{cor:ch42-finite-bijection}
Let $h:T\to X$ be a finite surjective morphism of normal integral
projective $k$-varieties and let $\theta\in N^1(X)_{\mathbf R}$.  Then
\[
 h^*:\PSH_X(\theta)\longrightarrow
 \bigl\{u\in\PSH_T(h^*\theta)\ \big|\
 u\ \text{is constant on the fibres of}\ h^{\An}\bigr\}
\]
is a bijection.
\end{corollary}

\begin{proof}
By \Cref{thm:bj-normal-envelope} applied to $T$ and $h^*\theta$, the
hypothesis $\operatorname{EP}_T(h^*\theta)$ of
\Cref{cor:ch41-finite-bijection} holds; that corollary is the assertion.
\end{proof}

\begin{remark}[These two corollaries are consequences, not inputs]
\label{rem:ch42-bijections-are-output}
\Cref{cor:ch42-birational-bijection} and \Cref{cor:ch42-finite-bijection}
are exactly the statements that \Cref{rem:ch39-one-way} and
\Cref{warn:ch41-one-way} declare unavailable inside
\Cref{chap:bj-reflection} and \Cref{chap:bj-finite-descent}.  They become
available here because \Cref{thm:bj-normal-envelope} has been proved, and
they are used nowhere in its proof.  A reader auditing the dependency
graph should check that neither corollary is cited above
\Cref{sec:ch42-consequences} in this chapter; it is not.
\end{remark}

\subsection{The hypothesis audit}
\label{sec:ch42-hypothesis-audit}

\begin{proposition}\label{prop:bj-alteration-hypothesis-audit}
Each hypothesis of \Cref{thm:bj-normal-envelope} has a distinct use:
\begin{enumerate}
\item integrality gives a single function field and an integral alteration;
\item projectivity makes the alteration source projective and places every
      analytic semivaluation in the compact BJ setting;
\item normality of $X$ makes finite normal descent applicable to $h$;
      normality of the Stein middle term follows from normality of $Y$;
\item algebraic closedness makes the ground field perfect, so the regular
      alteration source is smooth, and removes the finite ground-field
      extension in de Jong's theorem;
\item triviality of the absolute value is part of the BJ theory used in all
      three descent inputs;
\item the class is allowed to be real because each input theorem is stated for
      real numerical classes.
\end{enumerate}
No item in this list implies unibranchness of a later nonnormal target.
That additional bridge is the subject of
\Cref{chap:bj-universal-homeomorphisms,chap:bj-subintegral-normalization}.
\end{proposition}

\begin{proof}
The first five assertions were checked in the construction of
\eqref{eq:bj-alteration-stein} and in the applicability statements in the
proof of \Cref{thm:bj-normal-envelope}.
The sixth is the class identity
\eqref{eq:bj-alteration-class-ledger} together with the real-class statements
of the three input theorems.
The last sentence follows because the proof never replaces a nonnormal
variety by its normalization.

For the record, the individual verifications are these.  Integrality is used
in \Cref{lem:ch42-generic-finite} to identify the generic fibre with
$\Spec K(Y)$, hence in \eqref{eq:bj-alteration-function-fields}, and it is
the standing hypothesis under which $\PSH_X(\theta)$ is defined in
\Cref{def:bj-psh}.  Projectivity is used in the proof of
\Cref{thm:bj-de-jong-input} to see that the de Jong source is projective and
not merely quasi-projective, in \Cref{lem:bj-alteration-middle} for the
projectivity of $g$, and throughout \Cref{chap:bj-finite-descent} for the
flattening construction and for compactness of analytifications.
Normality of $X$ is used in Step~3 of
\Cref{thm:ch42-alteration-descent}, through
\Cref{lem:bj-finite-normal-open} and
\Cref{thm:bj-carrier-reflection}; normality of $T$ is used in the source psh/carrier input and in
birational descent along $g$, and is supplied by
\cite[\href{https://stacks.math.columbia.edu/tag/035L}{Tag 035L}]
{stacks-project}; normality of $Y$ is automatic, $Y$ being smooth.
The field hypotheses have two distinct geometric roles: perfectness makes a
regular finite-type source smooth, while algebraic closedness eliminates the
finite ground-field extension $k_1/k$.  Algebraic closedness is also part of
the valuation-theoretic standing conventions of \Cref{part:bj}.  Triviality
of the absolute value is
part of the definition of every object involved.  Real classes are allowed by the class-convergence and real-weight
interfaces of all three inputs.  Separately,
\Cref{lem:ch34-psh-nonempty} makes EP vacuous when the psh class
is empty; this is not a positivity hypothesis.
\end{proof}

\section{Four computations}
\label{sec:ch42-examples}

The four subsections that follow carry the chapter to explicit numbers on a
single surface.  \Cref{sec:ch42-cone-alteration} constructs a smooth
alteration of a singular normal surface and computes its Stein
factorization.  \Cref{sec:ch42-cone-endtoend} runs the descent of
\Cref{thm:ch42-alteration-descent} end to end on that surface, with an
explicit family of psh functions, through all five varieties that the proof
involves.  \Cref{sec:ch42-inseparable} changes the characteristic to two,
where the finite factor becomes purely inseparable and the trace map
vanishes identically, and checks that the determinant norm of
\Cref{sec:ch41-determinant} still supplies the required average.  \Cref{sec:ch42-no-separation} shows what fails when the
birational and the finite mechanisms are not separated.

The surface is the one already used in \Cref{chap:bj-finite-descent}, so that
the computations here and there can be compared.  We recall the notation.

\begin{itemize}
\item $X:=\mathbf P(1,1,2)$, the weighted projective plane with weighted
      coordinates $a,b$ of weight one and $z$ of weight two.  It is a
      normal integral projective surface, singular exactly at the vertex
      $o=[0:0:1]$.  Under
      $A=a^2,\ B=b^2,\ C=ab$ it is the quadric cone
      $\{AB=C^2\}\subset\mathbf P^3$, with $z$ the fourth coordinate and
      $o=[0:0:0:1]$.
\item $T:=\mathbf P^2$, with coordinates $[u:v:w]$, and
      \[
       h:T\longrightarrow X,\qquad [u:v:w]\longmapsto[u:v:w^2],
      \]
      a finite surjective morphism of degree two with
      $h^*a=u$, $h^*b=v$, $h^*z=w^2$.  When
      $\operatorname{char}k\ne2$ it is the quotient by the involution
      $\iota[u:v:w]=[u:v:-w]$.  We write
      $\tilde o:=[0:0:1]\in T$, the unique point over $o$.
\item $H$, the hyperplane class of $X\subset\mathbf P^3$, equal to
      $\cO_X(2)$ in weighted coordinates, with
      $H^0(X,H)$ spanned by $a^2,ab,b^2,z$.  It is a line bundle, whereas
      $\cO_X(1)$ is only a Weil divisor class and is not invertible at
      $o$.  We set $\theta:=c_1(H)\in N^1(X)_{\mathbf R}$.  Then
      \[
       h^*H=\cO_{\mathbf P^2}(2),\qquad
       (\theta^2)=2,\qquad
       (h^*\theta)^2=4=2\,(\theta^2),
      \]
      the middle equality because a quadric surface in $\mathbf P^3$ has
      degree two, and the last being
      \Cref{lem:ch42-class-ledger}\,\cref{itm:ch42-cl:degree} for $d=2$.
      Since $\operatorname{Cl}(X)=\mathbf Z\cdot\cO_X(1)$ and
      $\Pic(X)=\mathbf Z\cdot H$, the space $N^1(X)_{\mathbf R}$ is
      one-dimensional and is spanned by $\theta$.
\end{itemize}

\subsection{An alteration of a singular normal surface, with its Stein
factorization}
\label{sec:ch42-cone-alteration}

\begin{example}[A smooth alteration of the quadric cone and its Stein
factorization]
\label{ex:ch42-cone-stein}
Assume $\operatorname{char}k\ne2$ and keep the notation above.  Fix the
closed point
\[
 p:=[1:1:1]\in T=\mathbf P^2,
\]
and note that $p$ is not fixed by $\iota$: its $w$-coordinate is
nonzero and it is not $\tilde o$, whereas
$\operatorname{Fix}(\iota)=\{w=0\}\cup\{\tilde o\}$.  Write
$p':=\iota(p)=[1:1:-1]\ne p$.  Let
\[
 \sigma:Y:=\mathrm{Bl}_p\mathbf P^2\longrightarrow T
\]
be the blowup at $p$, with exceptional curve $E\cong\mathbf P^1$,
$E^2=-1$, and put
\[
 f:=h\circ\sigma:Y\longrightarrow X .
\]

\emph{Step 1: $f$ is a smooth projective alteration of degree two.}
The surface $Y$ is smooth, integral and projective, being a blowup of
$\mathbf P^2$ at a closed point; indeed $Y\cong\mathbf F_1$.  The
morphism $f$ is proper as a composite of proper morphisms, and surjective
because $\sigma$ and $h$ are.  It is generically finite of degree
\[
 \deg f=\deg h\cdot\deg\sigma=2\cdot1=2
\]
by \Cref{lem:ch42-degree-composite}\,\cref{itm:ch42-deg:comp}.  Thus $f$ is
an alteration in the sense of \Cref{def:bj-alteration}, with smooth
projective source; it is exactly the kind of morphism
\Cref{thm:bj-de-jong-input} produces abstractly, written down by hand.

\emph{Step 2: $f$ is not finite and not birational.}
It is not birational because $\deg f=2$.  It is not finite because it
contracts $E$: with $q:=h(p)=[1:1:1]\in X$ one has
\[
 f^{-1}(q)=\sigma^{-1}\bigl(h^{-1}(q)\bigr)
          =\sigma^{-1}(\{p,p'\})
          =E\ \sqcup\ \{\tilde p'\},
\]
where $\tilde p'$ denotes the unique point of $Y$ over $p'$, the
morphism $\sigma$ being an isomorphism near $p'\ne p$.  The fibre is
therefore the disjoint union of a curve and a point.  In particular
\Cref{prop:ch42-separation} is not vacuous for this $f$: both of its
factors are nontrivial, by
\Cref{prop:ch42-separation}\,\cref{itm:ch42-sep:degenerate}.

\emph{Step 3: the Stein factorization is $Y\to T\to X$.}
We claim $\sigma_*\cO_Y=\cO_T$.  The sheaf $\sigma_*\cO_Y$ is a coherent
sheaf of $\cO_T$-algebras, because $\sigma$ is proper; it is contained in
the constant sheaf $K(Y)=K(T)$, because $Y$ is integral and $\sigma$ is
birational; a coherent subsheaf of the function field is a finite
$\cO_T$-module, hence consists of elements integral over $\cO_T$; and
$T=\mathbf P^2$ is normal.  Therefore $\sigma_*\cO_Y=\cO_T$.

Consequently, by \Cref{lem:ch42-stein-universal}\,\cref{itm:ch42-su:unique}
applied with $X'=T$, $g'=\sigma$ and $h'=h$ --- the hypotheses being
that $h$ is finite, which it is, and that
$\sigma_*\cO_Y=\cO_T$, which has just been proved --- the Stein
factorization of $f$ is
\[
 Y\xrightarrow{\ \sigma\ }T\xrightarrow{\ h\ }X,
 \qquad
 g=\sigma,\qquad
 \Spec_X(f_*\cO_Y)=\Spec_X(h_*\cO_T)=T .
\]
The middle identification is the computation
$f_*\cO_Y=h_*\sigma_*\cO_Y=h_*\cO_T$.
It can also be seen through
\Cref{lem:ch42-middle-normalization}: the Stein middle term is the
normalization of $X$ in $K(Y)=K(\mathbf P^2)$, and $\mathbf P^2$ is
normal with that function field and finite over $X$, so it \emph{is} that
normalization.

Everything asserted by \Cref{prop:ch42-separation} is now visible:
$T=\mathbf P^2$ is normal integral projective with $K(T)=K(Y)$;
$g=\sigma$ is projective birational between normal integral projective
surfaces; and $h$ is finite surjective of degree $2=\deg f$ between
normal integral projective surfaces.

\emph{Step 4: the two fibres that matter.}
Over the vertex, $h^{-1}(o)=\{\tilde o\}$ as a set, since
$[u:v:w^2]=[0:0:1]$ forces $u=v=0$; and $\sigma$ is an isomorphism near
$\tilde o\ne p$, so $f^{-1}(o)$ is a single point.  The
\emph{scheme-theoretic} fibre of $h$ over $o$ is more interesting.  In
the chart $w\ne0$ upstairs, with $s=u/w$, $t=v/w$, the chart of $X$
is $R=k[s^2,st,t^2]$ with $\mathfrak m_o=(s^2,st,t^2)$, so
\[
 \cO_{h^{-1}(o)}=k[s,t]/(s^2,st,t^2),
 \qquad
 \dim_k=3,
\]
with $k$-basis $1,s,t$, whereas the generic rank of $h$ is
$d=2$.  The jump from $2$ to $3$ is precisely the failure of flatness
computed in \Cref{ex:ch41-nonflat-cone}.  Over $q$, by contrast, the fibre
of $h$ is the reduced scheme $\{p,p'\}$ of length two, because $\iota$
acts freely near $p$ and $\operatorname{char}k\ne2$, so $h$ is étale
there.

\emph{Step 5: the class ledger in numbers.}
The group $N^1(Y)_{\mathbf R}$ is two-dimensional, spanned by
$\sigma^*H_T$ and $E$, where $H_T:=c_1(\cO_{\mathbf P^2}(1))$; the
intersection numbers are
\[
 (\sigma^*H_T)^2=1,\qquad (\sigma^*H_T\cdot E)=0,\qquad E^2=-1 .
\]
With $\theta=c_1(H)$ one has
\[
 h^*\theta=2H_T,
 \qquad
 f^*\theta=\sigma^*h^*\theta=2\,\sigma^*H_T,
\]
which is \eqref{eq:bj-alteration-class-ledger} in this instance, and
\[
 (f^*\theta)^2=4(\sigma^*H_T)^2=4=2\,(\theta^2),
 \qquad
 (f^*\theta\cdot E)=2(\sigma^*H_T\cdot E)=0,
\]
the second identity recording that $E$ is contracted by $f$.  Note that
no factor $1/2$ appears anywhere: the birational factor $\sigma$
preserves intersection numbers and the finite factor $h$ doubles them,
exactly as \Cref{warn:ch42-no-degree} requires.

Finally, $f^*N^1(X)_{\mathbf R}=\mathbf R\cdot\sigma^*H_T$ is a line in the
plane $N^1(Y)_{\mathbf R}$; most classes upstairs, for instance
$\sigma^*H_T-E$, are not pullbacks at all.  This will be used in
\Cref{ex:ch42-no-separation}.
\end{example}

\begin{example}[The two degenerate cases]
\label{ex:ch42-stein-degenerate}
Keeping the notation above, two degenerations of
\Cref{ex:ch42-cone-stein} illustrate
\Cref{prop:ch42-separation}\,\cref{itm:ch42-sep:degenerate}.

\emph{The alteration is finite.}  Take $f:=h:T\to X$ itself, an alteration
of degree two with smooth projective source.  It is affine, so by
\Cref{lem:ch42-stein-universal}\,\cref{itm:ch42-su:finite} the birational
Stein factor is an isomorphism and the factorization is
$T\xrightarrow{\sim}T\xrightarrow{h}X$.  The descent of
\Cref{thm:ch42-alteration-descent} then consists of Step~3 alone, and it is
exactly \Cref{ex:ch41-descent-cone}.

\emph{The alteration is birational.}  Take $f:=\beta:Z=\mathbf F_2\to X$,
the minimal resolution of $X$ constructed in
\Cref{ex:ch41-flatten-cone}; it is a projective birational morphism with
smooth projective source.  Here $d=1$, so by
\Cref{lem:ch42-h-degree}\,\cref{itm:ch42-hd:h-iso} the finite Stein factor is
an isomorphism, $X$ being normal, and the factorization is
$Z\xrightarrow{\beta}X\xrightarrow{\sim}X$.  The descent consists of
Step~2 alone.  Since $\operatorname{char}k$ is arbitrary in this
particular instance --- the minimal resolution of a quadric cone exists in
every characteristic --- the birational descent alone supplies
$\operatorname{EP}_X(\theta)$ for this $X$.
It uses the route of \Cref{cor:ch42-char-zero-shortcut}, with an
explicit resolution available in every characteristic for this surface.
The argument uses the existence of a resolution in dimension two: the all-dimensional statement cannot be based on a
resolution theorem in positive characteristic, so
\Cref{thm:bj-de-jong-input} is indispensable there.
\end{example}

\begin{example}[A second factorization, with a different middle variety]
\label{ex:ch42-second-factorization}
Assume $\operatorname{char}k\ne2$ and keep the notation of
\Cref{sec:ch42-examples}.  This example exhibits an alteration of $X$ that
admits \emph{two} factorizations into a birational and a finite morphism ---
one in each order --- and computes both middle varieties.  It is the promised
illustration of \Cref{rem:ch42-universal-buys} and of
\Cref{warn:ch42-order}.

The alteration is not the one of \Cref{ex:ch42-cone-stein}.  Instead of
blowing up the general point $p=[1:1:1]$, blow up the point lying over the
\emph{vertex}:
\[
 \alpha:W:=\mathrm{Bl}_{\tilde o}\mathbf P^2\longrightarrow T=\mathbf P^2,
 \qquad
 \tilde o=[0:0:1]=h^{-1}(o)_{\mathrm{red}},
\]
and set $f':=h\circ\alpha:W\to X$.  As in
\Cref{ex:ch42-cone-stein}, $W\cong\mathbf F_1$ is smooth integral
projective, $f'$ is proper and surjective, and
$\deg f'=\deg h\cdot\deg\alpha=2$ by
\Cref{lem:ch42-degree-composite}\,\cref{itm:ch42-deg:comp}; so $f'$ is a
smooth projective alteration of $X$ of degree two, and
\Cref{thm:ch42-alteration-descent} applies to it.

\emph{The Stein factorization: birational, then finite.}
Exactly as in Step~3 of \Cref{ex:ch42-cone-stein}, $\alpha_*\cO_W=\cO_T$,
because $\alpha$ is proper birational with normal integral target.  Since
$h$ is finite,
\Cref{lem:ch42-stein-universal}\,\cref{itm:ch42-su:unique} identifies the
Stein factorization of $f'$ as
\[
 W\xrightarrow{\ \alpha\ }T=\mathbf P^2\xrightarrow{\ h\ }X,
 \qquad
 \Spec_X(f'_*\cO_W)=\Spec_X(h_*\cO_T)=\mathbf P^2 .
\]
The middle variety is $\mathbf P^2$, of Picard rank one.

\emph{A second factorization: finite, then birational.}
By \Cref{ex:ch41-flatten-cone}, the involution $\iota$ lifts to $W$ and
fixes the exceptional curve $F\subset W$ pointwise, the quotient
\[
 q:W\longrightarrow Z:=W/\iota\cong\mathbf F_2
\]
is finite flat of rank two between smooth projective surfaces, and the
contraction $\beta:Z\to X$ of the negative section $E=q(F)$, with
$E^2=-2$, is the minimal resolution of $X$.  The square of
\Cref{lem:bj-finite-flatification} commutes, that is,
\[
 \beta\circ q=h\circ\alpha=f' .
\]
So $f'$ also factors as a finite morphism of degree two followed by a
projective birational morphism, with middle variety $Z\cong\mathbf F_2$,
of Picard rank two.  Both middle varieties are normal integral projective
surfaces; they are not isomorphic, since $\mathbf F_2$ contains a curve of
negative self-intersection and $\mathbf P^2$ does not.

\emph{Which one is canonical.}  The Stein middle term is $\mathbf P^2$, and
$\mathbf F_2$ is not a Stein middle term of $f'$: the first factor of the
second factorization is $q$, and $q_*\cO_W$ is locally free of rank two,
not $\cO_Z$, so the hypothesis of
\Cref{lem:ch42-stein-universal}\,\cref{itm:ch42-su:unique} fails for it.
Equivalently, by \Cref{lem:ch42-middle-normalization} the Stein middle term
is the normalization of $X$ in $K(W)=K(\mathbf P^2)$, which is
$\mathbf P^2$ itself; the surface $\mathbf F_2$ has function field
$K(X)$ and is therefore a birational model of $X$, not a cover of it.

\emph{What each factorization does with the exceptional curve.}  The two
factorizations dispose of the curve $F\subset W$ at opposite ends.  In the
Stein factorization $F$ is contracted immediately, by $\alpha$, and the
finite morphism $h$ is then applied on $\mathbf P^2$, where the vertex
has been re-created.  In the second factorization $F$ survives the finite
morphism --- $q$ maps it isomorphically onto $E$ --- and is contracted
only at the last step, by $\beta$.  This is the concrete content of the
asymmetry recorded in \Cref{rem:ch42-universal-buys}: the universal property
produces the largest affine factor \emph{under} $W$, and there is no
universal property producing a smallest finite factor \emph{over} $X$.

\emph{Both orders happen to work here, and this proves nothing in general.}
In this instance one may run the descent either way.  Along the Stein
factorization: $\operatorname{EP}_W(f'^*\theta)$ from
\Cref{thm:bj-smooth-envelope}, then
\Cref{thm:bj-birational-ep-descent} along $\alpha$, then
\Cref{thm:bj-finite-normal-descent} along $h$.  Along the second
factorization: $\operatorname{EP}_W(f'^*\theta)$ again, then
\Cref{thm:bj-finite-normal-descent} along $q$ --- legitimate, since
$W$ and $Z$ are smooth, hence normal, integral projective and $q$ is
finite surjective --- and then \Cref{thm:bj-birational-ep-descent} along
$\beta$.  Both routes end at $\operatorname{EP}_X(\theta)$.

This reverse-order factorization is a special feature of the explicit
surface computation.  Flattening the finite map $h$ produces the square,
and the original source $W$ was chosen to be its upper-left term.
For a general alteration $Y\to X$, flattening its Stein factor
need not give a map from that unchanged $Y$ to the modified target.
Here normality and smoothness of $Z$ follow from explicit blowup
charts, rather than from a general assertion about smooth fixed loci.
If a reverse factorization has a nonnormal middle term, the
normalization argument in \Cref{warn:ch42-order} repairs it.
\end{example}

\subsection{The descent, run end to end}
\label{sec:ch42-cone-endtoend}

We now run the proof of \Cref{thm:ch42-alteration-descent} on the surface
$X=\mathbf P(1,1,2)$ with an explicit bounded-above family of psh
functions, computing every intermediate object in closed form.  Five
varieties occur: the smooth alteration source $Y$, the Stein middle term
$T$, the base $X$, and, inside the finite step, the two terms $W$ and
$Z$ of the flattening square.  The family has been chosen so that all five
computations can be carried out by hand; the price is that in this instance
the supremum happens to be upper semicontinuous already, so the
regularization $(\,\cdot\,)^*$ is not exercised.  Where regularization does
bite is shown in \Cref{ex:ch39-two-usc-extensions} and
\Cref{ex:ch41-node-not-open}.

\begin{example}[The descent of \Cref{thm:ch42-alteration-descent}, computed]
\label{ex:ch42-endtoend}
Assume $\operatorname{char}k\ne2$, keep the notation of
\Cref{sec:ch42-examples}, and let
\[
 f=h\circ\sigma:Y=\mathrm{Bl}_p\mathbf P^2\longrightarrow X,
 \qquad p=[1:1:1],
\]
be the smooth projective alteration of \Cref{ex:ch42-cone-stein}, with Stein
factorization $Y\xrightarrow{\sigma}T=\mathbf P^2\xrightarrow{h}X$ and
$\theta=c_1(H)$, $H=\cO_X(2)$.

\emph{Stage 0: the family and its supremum.}
For $c\in k^\times$ put
\[
 \varphi_c:=\log\bigl|a^2-cz\bigr|\in\cH^{\mathrm{gf}}(H)\subset
 \PSH_X(\theta),
\]
which is legitimate because $a^2-cz$ is a nonzero global section of the
line bundle $H$.  The family $(\varphi_c)_{c\in k^\times}$ is nonempty,
and it is bounded above by $0$: indeed
$\sup_{X^{\An}}\log|s|=\log|s|(v_{X,\triv})=0$ for every nonzero
$s\in H^0(X,H)$, by
\Cref{lem:ch34-log-s-properties}\,\cref{itm:ch34-logs:max} applied with the
irreducible subvariety $X$ itself.  It is therefore an admissible family
for $\operatorname{EP}_X(\theta)$.

We claim that, pointwise on $X^{\An}$,
\[
 \sup_{c\in k^\times}\varphi_c
 =\max\bigl\{\log|a^2|,\ \log|z|\bigr\}=:u_X ,
\]
with no regularization needed.  The inequality $\le$ is
\Cref{lem:ch34-log-s-properties}\,\cref{itm:ch34-logs:sum} together with
$\log|cz|=\log|z|$, the absolute value being trivial on $k$.  For the
reverse inequality fix $x\in X^{\An}$.  If $z$ vanishes identically on
the support of $x$, then $\log|a^2-cz|(x)=\log|a^2|(x)$ for every $c$
and both sides agree; the same argument applies if $a^2$ does.  Otherwise
both values are finite and we may form
$\zeta:=(a^2/z)(x)$ in the completed residue field $\mathcal H(x)$, so
that
\[
 \varphi_c(x)=\log|z|(x)+\log|\zeta-c| .
\]
If $|\zeta|\ne1$ then $|\zeta-c|=\max\{|\zeta|,1\}$ for every
$c\in k^\times$ and we are done.  If $|\zeta|=1$, then
$|\zeta-c|<1$ forces $c$ to be the image in the residue field of
$\mathcal H(x)$ of the residue of $\zeta$; since the absolute value is
trivial on $k$, the composite $k\hookrightarrow\mathcal H(x)^{\circ}\to
\widetilde{\mathcal H(x)}$ is injective, so at most one $c\in k^\times$
can fail, and $k$ is infinite.  This proves the claim.

Since $u_X$ is again a Fubini--Study function --- the maximum of
$\log|a^2|$ and $\log|z|$ with equal weights, both sections of $H$ ---
it is upper semicontinuous and lies in $\PSH_X(\theta)$.  That is the
conclusion $\operatorname{EP}_X(\theta)$ asks for, for this family; the
point of what follows is to obtain it by the route of
\Cref{thm:ch42-alteration-descent} rather than by inspection.

\emph{Stage 1: upstairs on the smooth source.}
By \Cref{lem:ch41-proper-gf-pullback}\,\cref{itm:ch41-pgf:values} the
pullbacks are computed section by section.  Since $h^*a=u$, $h^*b=v$,
$h^*z=w^2$, we get $h^*(a^2-cz)=u^2-cw^2$, and hence
\[
 f^*\varphi_c=\log\bigl|u^2-cw^2\bigr|
 \in\cH^{\mathrm{gf}}(f^*H)\subset\PSH_Y(f^*\theta),
 \qquad f^*\theta=2\,\sigma^*H_T .
\]
The identical ultrametric argument on $Y^{\An}$, with $u^2/w^2$ in place
of $a^2/z$, gives
\[
 \sup_{c}f^*\varphi_c
 =\max\bigl\{\log|u^2|,\log|w^2|\bigr\}\circ\sigma^{\An}
 =f^*u_X=:u_Y ,
\]
again with no regularization.  The membership $u_Y\in\PSH_Y(f^*\theta)$ is
what \Cref{thm:bj-smooth-envelope} guarantees abstractly, $Y$ being a
smooth projective variety and $f^*\theta$ an arbitrary real class; here one
also sees it directly, $u_Y$ being visibly Fubini--Study.  This illustrates how the smooth theorem supplies the initial EP.
For this family its supremum is already an explicit FS function,
so membership can also be checked without that theorem.

\emph{Stage 2: the birational descent, from $Y$ to $T$.}
Put
\[
 u_T:=\max\bigl\{2\log|u|,\ 2\log|w|\bigr\}
 \in\cH^{\mathrm{gf}}\bigl(\cO_{\mathbf P^2}(2)\bigr)
 \subset\PSH_T(h^*\theta),
 \qquad h^*\theta=2H_T .
\]
Then $\sigma^*u_T=u_Y$, by
\Cref{lem:ch41-proper-gf-pullback}\,\cref{itm:ch41-pgf:values} again.  In the
proof of \Cref{thm:ch42-alteration-descent} this stage is not a computation
with one function but the transfer of a \emph{property}: Step~2 applies
\Cref{thm:bj-birational-ep-descent} to $\sigma$ and concludes
$\operatorname{EP}_T(h^*\theta)$.  What the present family shows is the
shape of the objects involved: the family upstairs is
$\sigma^*(h^*\varphi_c)$, its supremum is a pullback, and the reflection
returns $u_T$.  The classes match without a degree factor, as
\Cref{warn:ch42-no-degree} predicts:
$\sigma^*(2H_T)=2\sigma^*H_T=f^*\theta$, and
$(f^*\theta)^2=4=2\,(\theta^2)$, the factor two being contributed by
$h$ and not by $\sigma$.

\emph{Stage 3: the finite descent, from $T$ to $X$.}
Step~0 of \Cref{thm:bj-finite-normal-descent} forms
$(\sup_ch^*\varphi_c)^*$ and identifies it with $h^*u_X$, the
identification being the openness statement of
\Cref{lem:bj-finite-normal-open}.  Here $h^*u_X=u_T$, as computed.

Step~1 flattens.  By \Cref{ex:ch41-flatten-cone} the flattening square for $h$ is
\[
 W=\mathbf F_1\xrightarrow{\ \alpha\ }T=\mathbf P^2,
 \qquad
 q:W\to Z=\mathbf F_2,
 \qquad
 \beta:Z\to X,
\]
with $q$ finite flat of rank two and $\beta$ the minimal resolution.
Note that $W$ and $Y$ are both isomorphic to $\mathbf F_1$ and are
\emph{different} surfaces over $T$: $Y$ is the blowup at
$p=[1:1:1]$, while $W$ is the blowup at $\tilde o=[0:0:1]$.  Confusing
them is the error that \Cref{warn:ch42-stein-not-flattening} warns against.

Now write $u_T=\max_j\{\log|s_j|\}$ with $m=1$,
$M:=h^*H=\cO_{\mathbf P^2}(2)$, $s_0:=u^2$, $s_1:=w^2$ and weights
$\lambda_0=\lambda_1=0$.  Because $s_0=h^*(a^2)$ and $s_1=h^*(z)$, and
because $\alpha^*h^*=q^*\beta^*$, the pulled-back sections are themselves
pullbacks along $q$:
\[
 \alpha^*s_0=q^*\bigl(\beta^*a^2\bigr),
 \qquad
 \alpha^*s_1=q^*\bigl(\beta^*z\bigr),
 \qquad
 \alpha^*M=q^*\bigl(\beta^*H\bigr).
\]
The universal norm polynomial of \Cref{lem:ch41-universal-polynomial} is
therefore computed by the multiplicativity of the norm on pullbacks,
$\mathrm N_q(q^*y)=y^{\otimes2}$:
\[
 \begin{aligned}
 \mathrm N_q\bigl(T_0\,\alpha^*s_0+T_1\,\alpha^*s_1\bigr)
 &=\bigl(T_0\,\beta^*a^2+T_1\,\beta^*z\bigr)^{2}\\
 &=T_0^2\,(\beta^*a^2)^2+2\,T_0T_1\,\beta^*(a^2z)+T_1^2\,(\beta^*z)^2 .
 \end{aligned}
\]
Writing $A:=\log|\beta^*a^2|$ and $B:=\log|\beta^*z|$, the definition of
the averaging operator, with $m=1$ and $d=2$, gives
\[
 \operatorname{Av}_q(\alpha^*u_T)
 =\tfrac12\max\bigl\{2A,\ A+B,\ 2B\bigr\}
 =\max\{A,B\}
 =\beta^*u_X ,
\]
the middle term being harmless because $A+B\le2\max\{A,B\}$ and
$\log|2|=0$.  The classes match:
$\tfrac12q_*\alpha^*(2H_T)=\tfrac12q_*q^*\beta^*\theta=\beta^*\theta$,
using $q_*q^*=2$ from \eqref{eq:ch41-projection-formula}, and
$\beta^*u_X$ has class $\beta^*\theta$.

The carrier condition in Step~2 is explicit here.
Since $\cO_T(2)$ is globally generated, $0$ is a homogeneous FS
function in its class.  The normalized carrier family contains zero
and every member is at most zero, so $P_T=0$.
Its average is $p_Z=0\in\PSH_{\mathrm{hom},Z}(\beta^*\theta)$,
by the norm limit formula or global generation of $\beta^*H$.
The computed function satisfies $\beta^*u_X\le p_Z$.

Step~3 reflects $\beta^*u_X\in\PSH_Z(\beta^*\theta)$ down the projective
birational morphism $\beta$ by \Cref{thm:bj-carrier-reflection}, producing
a function on $X$ whose pullback is $\beta^*u_X$; since $\beta^{\An}$
is surjective, that function is $u_X$.  Hence
$u_X\in\PSH_X(\theta)$, which is the conclusion.

\emph{The fibre formula, checked at two points.}
The identity \eqref{eq:bj-finite-flat-average} can be verified numerically.
Take first the smooth point $P:=[0:1:0]\in X$, and let $v_P$ be the
divisorial valuation given by the exceptional divisor of the blowup of $X$
at $P$.  In the affine chart $b\ne0$ of $X$, with coordinates
$x=a/b$ and $y=z/b^2$ and frame $b^2$ of $H$, one has $v_P(x)=v_P(y)=1$,
whence
\[
 \log|a^2|(v_P)=-2,\qquad \log|z|(v_P)=-1,\qquad u_X(v_P)=-1 .
\]
Upstairs, in the chart $v\ne0$ of $T$ with coordinates $u,w$, the
morphism reads $h^*x=u$, $h^*y=w^2$, so the unique point of $T^{\An}$
over $v_P$ is the monomial valuation $w_P$ with
$w_P(u)=1$ and $w_P(w)=\tfrac12$; the fibre algebra over $v_P$ is
$\mathcal H(v_P)[w]/(w^2-y)$, a field, since $|y|=\rme^{-1}$ is not a
square in the value group $\rme^{\mathbf Z}$.  Hence there is one point in
the fibre, with weight $e(w_P/v_P)=2$, and
\[
 \begin{aligned}
  u_T(w_P)&=\max\{-2,-1\}=-1,\\
  \tfrac12\sum_{w\mid v_P}e(w/v_P)\,u_T(w)
   &=\tfrac12\cdot2\cdot(-1)=-1=u_X(v_P).
 \end{aligned}
\]
Take next the divisorial valuation $v_E:=\operatorname{ord}_E$ attached to
the negative section $E\subset Z$.  Near the vertex the nonvanishing
section $z$ frames $H$, and $a^2/z$ is the coordinate $s^2$ in the
resolution chart $k[s^2,\rho]$ of
\Cref{ex:ch41-flatten-cone}.  The divisor $E$ is cut out by $s^2$,
so the coefficient of $\beta^*a^2$ has order one and that of
$\beta^*z$ is the unit $1$ with order zero.
Thus $\beta^*u_X(v_E)=\max\{-1,0\}=0$.
The numerical identity $\beta^*H=2F'+E$ alone would not determine
these individual section orders.
Its unique preimage under $q^{\An}$ is $\tfrac12\operatorname{ord}_F$,
because $q^*E=2F$, again with weight two, and
$\alpha^*u_T(\tfrac12\operatorname{ord}_F)=\max\{-1,0\}=0$; the average is
$\tfrac12\cdot2\cdot0=0$, as required.

\emph{What the computation illustrates and what it does not.}
It illustrates the class ledger, the two weights, the absence of a degree
factor in the birational step, and the distinctness of the five varieties.
It does not illustrate the regularization: here
$\sup_c\varphi_c$ is already upper semicontinuous, so
$\bigl(\sup_c\varphi_c\bigr)^*=\sup_c\varphi_c$ and the whole content of
$\operatorname{EP}$ is invisible.  The carrier is the zero carrier just checked, so no nontrivial carrier
tail is needed.  The functions themselves are not bounded on both sides:
$u_X(v_{P,\triv})=-\infty$ at $P=[0:1:0]$, the common zero of
$a^2$ and $z$, and the upstairs pullbacks have the corresponding poles.
Being FS does not imply boundedness.  The example illustrates domination
by a bounded carrier, leaving the general tail and regularization
mechanisms to \Cref{chap:bj-reflection}.
\end{example}

\subsection{An inseparable finite part}
\label{sec:ch42-inseparable}

The finite Stein factor of an alteration can be purely inseparable, and then
every construction that averages by summing over sheets degenerates: the
sheets are not there, the trace vanishes identically, and there is no group
to average over.  The determinant norm of \Cref{sec:ch41-determinant}
survives untouched, because it is a determinant and not a sum, and because
the number $d$ it divides by is a real normalization and not an element of
$k$.  The following example is the simplest one in which this can be
watched happening.

\begin{example}[The quadric cone in characteristic two]
\label{ex:ch42-inseparable-cone}
Let $\operatorname{char}k=2$ and keep
$X=\mathbf P(1,1,2)$, $T=\mathbf P^2$, and
\[
 h:T\longrightarrow X,\qquad [u:v:w]\longmapsto[u:v:w^2].
\]
Nothing in the definition of $h$ referred to the characteristic; what
changes is everything about its fibres.

\emph{(a) The extension of function fields is purely inseparable of degree
two.}
In the affine charts of Step~4 of \Cref{ex:ch42-cone-stein}, with
$s=u/w$ and $t=v/w$, the morphism is the inclusion of the second
Veronese subring,
\[
 R=k[s^2,st,t^2]\subset k[s,t],
\]
so $K(X)=k(s^2,st)$ and $K(T)=k(s,t)$.  Since
$(st)^2=s^2t^2$, we have $t^2=(st)^2/s^2\in K(X)$, whence
$K(X)\supseteq k(s^2,t^2)=K(T)^2$; therefore every element of $K(T)$ has
its square in $K(X)$, and $K(T)/K(X)$ is purely inseparable.  Its degree
is two: $[K(T):k(s^2,t^2)]=4$ because $s$ and $t$ are independent
$p$-th roots, while $[K(X):k(s^2,t^2)]=2$ because $(st)^2\in k(s^2,t^2)$
and $st\notin K(T)^2=k(s^2,t^2)$, the element $st$ having odd degree in
$s$.  Thus $d=\deg h=2=\operatorname{char}k$.

This is the situation announced in
\Cref{warn:ch42-de-jong-scope}\,(2): the alteration
$f=h\circ\sigma$ of \Cref{ex:ch42-cone-stein}, formed with any blowup
$\sigma$ of $\mathbf P^2$ at a closed point, is a smooth projective
alteration of $X$ of degree two whose function-field extension is
inseparable.  Its Stein factorization is unchanged, by the same computation
$\sigma_*\cO_Y=\cO_T$, and
\Cref{thm:ch42-alteration-descent} applies to it verbatim.

\emph{(b) The analytification of $h$ is a homeomorphism.}
Because $T$ and $X$ are normal integral projective and $K(T)/K(X)$ is
purely inseparable, $h$ is radicial, and
$h^{\An}:T^{\An}\to X^{\An}$ is a homeomorphism by
\Cref{lem:ch41-radicial-homeo}.  Every fibre of $h^{\An}$ is a single point.
The map $h$ is nevertheless nonflat at the vertex, where its scheme
fibre has length three.  The weights of
\Cref{lem:ch41-fibre-decomposition} apply to the finite flat $q$
constructed in part~(c), whose unique analytic lift has weight two.

The contrast with $\operatorname{char}k\ne2$ is sharpest at the point
$Q=[1:0:1]\in X$ used in \Cref{ex:ch42-not-a-pullback}.  There the fibre
of $h$ consists of the two reduced points $[1:0:1]$ and $[1:0:-1]$,
each of weight one.  In characteristic two the equation $w^2=1$ has the
single root $w=1$ with multiplicity two, and its scheme fibre is
$\Spec k[w]/(w-1)^2$.  This lies in the flat locus, where fibre
length and the norm weight are indeed two.  In particular \Cref{ex:ch42-not-a-pullback} has no analogue
here: since $h^{\An}$ is a homeomorphism, \emph{every} function on
$T^{\An}$ is the transport of a function on $X^{\An}$, and the
obstruction that example exhibits is a separable phenomenon.

\emph{(c) The flattening square is the same, and $q$ is again finite flat
of rank two.}
The local computation of \Cref{ex:ch41-flatten-cone} did not use the
characteristic.  Blow up $\tilde o=[0:0:1]$ upstairs and the vertex $o$
downstairs.  In the chart $t=s\rho$ one has
\[
 \cO_W=k[s,\rho],
 \qquad
 \cO_Z=k\bigl[s^2,\ st/s^2,\ t^2/s^2\bigr]=k[s^2,\rho,\rho^2]=k[s^2,\rho],
\]
and the other chart is obtained by exchanging $s$ and $t$.  These two affine planes cover the preimage of the vertex chart.
Off the vertex $\beta$ is an isomorphism and $X$ is smooth, so
$Z$ is smooth everywhere.  The characteristic-independent blowup
construction identifies it with the minimal resolution $\mathbf F_2$.
Moreover,
\[
 q:W=\mathbf F_1\longrightarrow Z=\mathbf F_2,
 \qquad
 k[s^2,\rho]\subset k[s,\rho]\ \text{free with basis}\ \{1,s\},
\]
is free of rank two in these charts.  On the remaining charts
$a\ne0$ and $b\ne0$, the map adjoins $w/u$ or $w/v$ with its
square prescribed, giving free bases $1,w/u$ or $1,w/v$.
Thus $q$ is finite locally free of rank two globally.
In characteristic two the sign automorphism becomes the identity.
There is no nontrivial constant order-two group giving this purely
inseparable quotient; this does not rule out nonreduced group schemes.  The identity $q^*E=2F$ persists --- $E$ is cut by $s^2$ and
$F$ by $s$ --- so $E^2=-2$ exactly as before, and the weight two of
part~(b) is visible as $\operatorname{ord}_F=2\operatorname{ord}_E$ on
$K(Z)=K(X)$.

\emph{(d) The trace vanishes identically.}
On the free rank-two algebra $k[s,\rho]\supset k[s^2,\rho]$ with basis
$\{1,s\}$, multiplication by $\xi=A+Bs$, $A,B\in k[s^2,\rho]$, has
matrix
\[
 \begin{pmatrix}A & Bs^2\\ B& A\end{pmatrix},
 \qquad
 \operatorname{Tr}(\xi)=2A=0,
 \qquad
 \det(\xi)=A^2-B^2s^2=(A+Bs)^2 .
\]
So the trace form is the zero map.  Any attempted average of the shape
``$\frac1d\log|\operatorname{Tr}(\cdot)|$'' returns $-\infty$
identically, on every function and at every point; it does not merely lose
constants, it loses everything.  Nor is there a Reynolds operator to fall
back on: the section operator $\frac12(1+\iota^*)$ available when
$\operatorname{char}k\ne2$ is unavailable here, since $2=0$ and
the sign automorphism is trivial.  Even in the separable case this
section operator does not compute the required logarithmic FS average.  This is the promised failure mode of
\Cref{rem:ch41-norm-vs-trace}, in its most extreme form; in
characteristic zero the trace already fails, but only by an inequality,
whereas here it fails by a collapse.

\emph{(e) The universal norm polynomial survives the degeneration.}\par\noindent
Let $\varphi=\max\{\log|s_0|+\lambda_0,\ \log|s_1|+\lambda_1\}$ be a
Fubini--Study function on $W^{\An}$ presented with $m=1$ and two sections
of a line bundle $M$, with $\lambda_0,\lambda_1\in\mathbf R$, which is
permitted by \Cref{cor:ch41-real-weights}.  Since $\det\xi=\xi^2$ by
part~(d), shrink an open set $U\subset Z$ so that $M$ is trivial on
$q^{-1}(U)$, choose a frame $e$ there, and write $s_i=b_i e$.  In the
induced frame $\mathrm N_q(e)$ of
the norm line, the universal norm polynomial of
\Cref{lem:ch41-universal-polynomial} is
\[
 \mathrm N_q\bigl(T_0s_0+T_1s_1\bigr)
 =\bigl(T_0b_0+T_1b_1\bigr)^2\mathrm N_q(e)
 =\bigl(T_0^2b_0^2+T_1^2b_1^2\bigr)\mathrm N_q(e),
\]
the cross term having vanished.  Thus, invariantly,
\[
 c_{(2,0)}=\mathrm N_q(s_0),\qquad c_{(1,1)}=0,\qquad
 c_{(0,2)}=\mathrm N_q(s_1).
\]
The mixed coefficient vanishes identically.  This is exactly the
degeneration that a proof organized around ``the $d$ elementary symmetric
functions of distinct sheets'' would miss, and it is harmless here: the
maximum defining the averaged function simply omits the vanishing
coefficient.  At $z\in Z^{\An}$, with unique preimage $w$, the norm formula
gives
\[
 \begin{aligned}
 \operatorname{Av}_q(\varphi)(z)
 &=\tfrac12\max\bigl\{
   \log|\mathrm N_q(s_0)|(z)+2\lambda_0,\ -\infty,\
   \log|\mathrm N_q(s_1)|(z)+2\lambda_1\bigr\}
 \\
 &=\max\bigl\{\log|s_0|(w)+\lambda_0,\
               \log|s_1|(w)+\lambda_1\bigr\}
 =\varphi(w).
 \end{aligned}
\]
This remains valid when one or both pure coefficients vanish: a vanishing
term contributes $-\infty$, and if both vanish then both sides are
$-\infty$.  It is also exactly the fibre formula
\eqref{eq:bj-finite-flat-average}: there is one point over each $z$, its
weight is two, and
$\frac12\cdot2\cdot\varphi(w)=\varphi(w)$.

\emph{(f) A closed-form instance, identical to the separable one.}
Take the section $u-w$ of $\cO_{\mathbf P^2}(1)$ of
\Cref{ex:ch42-not-a-pullback} and $\varphi'':=\log|u-w|$.  In the
chart above, with frame $w$, the element to be normed is $s-1$, and
\[
 \mathrm N_q(s-1)
 =\det\begin{pmatrix}-1& s^2\\ 1& -1\end{pmatrix}
 =1-s^2 .
\]
For the global section identity, the divisor argument of
\Cref{ex:ch41-descent-cone} applies in this characteristic as well:
the strict transform of $\{w=0\}$ maps with residue degree one onto
$\{z=0\}$, which avoids the vertex.  The norm divisor identity
identifies $\mathrm N_q(\alpha^*\cO(1))$ with $\beta^*H$ and sends
$\mathrm N_q(\alpha^*w)$ to $\beta^*z$.
The coefficient $1-s^2=(z-a^2)/z$ therefore sends the norm section
to $\beta^*(z-a^2)$.  Hence
\[
 \operatorname{Av}_q(\alpha^*\varphi'')=\tfrac12\,\beta^*\log\bigl|a^2-z\bigr| ,
\]
the same closed form as in \Cref{ex:ch41-descent-cone}.  The two
computations of the determinant are different, however: when
$\operatorname{char}k\ne2$ it is the product $(1-s)(1+s)$ over two
distinct sheets, and when $\operatorname{char}k=2$ it is the square
$(1+s)^2$ of a single sheet counted twice.  Only the second is available
here, and the determinant supplies its multiplicity.

\emph{(g) The two divisions by $d$ are not divisions in $k$.}
Here $\operatorname{char}k\mid d$, so a proof that divided by $d$
anywhere inside the algebra would collapse.  It does not: the factor
$1/(md)$ in the definition of $\operatorname{Av}_q$ divides a real
number, and the factor $1/d$ in
\eqref{eq:bj-norm-class} divides a numerical class in
$N^1(Z)_{\mathbf R}$.  Both are divisions in $\mathbf R$, and both are
legitimate for every $d\ge1$ in every characteristic; this is the content
of \Cref{warn:ch42-no-degree} and \Cref{rem:ch41-division-by-d}.

\emph{(h) What this example does and does not establish.}
It establishes that the finite factor of a perfectly ordinary alteration of
a perfectly ordinary normal surface can be purely inseparable, so that the
generality claimed for \Cref{chap:bj-finite-descent} is used and not
decorative.  It does not establish that the envelope property for this
particular $X$ requires the machinery: the minimal resolution
$\beta:\mathbf F_2\to X$ exists in every characteristic, so
\Cref{ex:ch42-stein-degenerate} already gives
$\operatorname{EP}_X(\theta)$ by a single birational descent.  What has no
resolution-based replacement in arbitrary dimension and characteristic is
the general mechanism; the example shows that mechanism running in a place
where it can be checked.  Finally, the fact that $h$ is here a universal
homeomorphism is a special feature of this example.  The proof of the general
normal case uses \Cref{thm:bj-finite-normal-descent}, whose inseparable
openness reduction already uses \Cref{lem:ch41-radicial-homeo}.  Their systematic treatment is the subject
of \Cref{chap:bj-universal-homeomorphisms}, which is not available in this
chapter and is not used.
\end{example}

\subsection{What fails without the separation}
\label{sec:ch42-no-separation}

The three defects of an alteration listed in
\Cref{sec:ch42-why-alteration} --- the loss of the common function field,
the non-finiteness of the fibres, and the fact that neither defect repairs
the other --- were stated there as observations about morphisms.  We now
carry each of them to an explicit psh function on an explicit surface.

\begin{example}[Three ways of trying to descend along the alteration itself]
\label{ex:ch42-no-separation}
Assume $\operatorname{char}k\ne2$ and keep the alteration
$f=h\circ\sigma:Y=\mathrm{Bl}_p\mathbf P^2\to X=\mathbf P(1,1,2)$ of
\Cref{ex:ch42-cone-stein}, with $p=[1:1:1]$, exceptional curve
$E_\sigma\subset Y$ over $p$, and $q_0:=h(p)=[1:1:1]\in X$.

\emph{(1) There is no common function field, and the reflection machinery
has nothing to act on.}
Here $K(Y)=K(\mathbf P^2)$ is a degree-two extension of $K(X)$.  The
whole of \Cref{chap:bj-reflection} is built on the identification of the
divisorial loci $X^{\divis}$ and $Y^{\divis}$ provided by
\Cref{lem:ch39-beta-an}\,\cref{itm:ch39-ban:val}, and that identification is
simply false here.  A single divisorial valuation of $K(X)$ can have
several extensions, or one extension with a weight: at the point
$Q=[1:0:1]\in X$ the trivial semivaluation has the two preimages
$v_{p_+,\triv}$ and $v_{p_-,\triv}$ of \Cref{ex:ch42-not-a-pullback},
while at the divisorial valuation $\operatorname{ord}_{C_a}$ attached to
the ruling $C_a=\{a=0\}$ there is a single preimage, the extension
carrying a residue field extension of degree two, because
$h|:\{u=0\}\to C_a$ is itself a degree-two map of projective lines.
The closed-point pair lies in proper support, whereas the ruling
valuation is divisorial.  Both show why the birational common-field
identification cannot be used for this cover.
The function $\sigma^*\log|u-w|$ is already nonconstant on the two
closed lifts of $Q$ and cannot be a pullback of any function on
$X^{\An}$.  This is the precise function-level obstruction.
It does not assert that every homogeneous psh function is a Cartier
coefficient function, or that one nonsplit divisorial lift alone
obstructs every possible descent construction.

\emph{(2) The fibres are not finite, and psh functions are not constant on
them.}
By Step~2 of \Cref{ex:ch42-cone-stein},
\[
 f^{-1}(q_0)=E_\sigma\ \sqcup\ \{\tilde p'\},
\]
the disjoint union of a rational curve and a point.  A weighted average of
the values of a function on this fibre, of the shape
\eqref{eq:ch41-average-shape}, is not defined: the sum has infinitely many
terms, and no system of positive integers summing to $d=2$ can be attached
to the points of a curve.

One might hope that this does not matter, on the ground that the functions
one has to average are constant on $E_\sigma$ anyway.  For
\emph{homogeneous} psh functions whose class is a pullback along $\sigma$
that is true, and it is \Cref{cor:ch39-fibre-constant}.  But it is not true
for the objects the machinery actually manipulates, because their classes
are not pullbacks.  Recall from \Cref{def:bj-psh} that a psh function is a
decreasing limit of Fubini--Study functions whose classes \emph{converge} to
the given class without being equal to it; and, as recorded at the end of
\Cref{ex:ch42-cone-stein}, $f^*N^1(X)_{\mathbf R}=\mathbf R\cdot\sigma^*H_T$
is a line inside the plane $N^1(Y)_{\mathbf R}$, so the approximating
classes are allowed to lie off it.

Here is a sequence of individually admissible FS stages with that
class convergence.  It is not a decreasing regularization.  For $m\ge1$
let
\[
 M_m:=\sigma^*\cO_{\mathbf P^2}(2m)\otimes\cO_Y(-E_\sigma).
\]
Its sections are degree-$2m$ forms vanishing at $p$.  Consider
\[
 P_{1,m}:=(u-w)\,v\,u^{2m-2},
 \qquad
 P_{2,m}:=(v-w)\,u^{2m-1}.
\]
Both vanish at $p=[1:1:1]$ to order exactly one; let $s_{1,m},s_{2,m}$ be
their strict transforms, viewed as sections of $M_m$.  Their
restrictions to $E_\sigma\cong\mathbf P(T_p\mathbf P^2)$ vanish at the
points $e_1,e_2$ determined by the tangent directions of $\{u=w\}$ and
$\{v=w\}$ at $p$, and $e_1\ne e_2$ because the two lines are distinct.
Put
\[
 \psi_m:=\frac1m\max\bigl\{\log|s_{1,m}|,\
                              \log|s_{2,m}|-1\bigr\}.
\]
This is a generically finite Fubini--Study function of class
\[
 \frac1m c_1(M_m)=2\sigma^*H_T-\frac1m[E_\sigma]
 \longrightarrow 2\sigma^*H_T=f^*\theta.
\]
At the trivial semivaluation of a closed point $e\in E_\sigma$ with
$e\ne e_1$ one has $s_{1,m}(e)\ne0$ and $\psi_m(e)=0$, while at
$e_1$ one has $s_{1,m}(e_1)=0$, $s_{2,m}(e_1)\ne0$, and
$\psi_m(e_1)=-1/m$.  Thus every $\psi_m$ is nonconstant on the
positive-dimensional fibre $f^{-1}(q_0)$ even though its presented
class tends to $f^*\theta$.  At $e_1$ the values $-1/m$ increase,
so this sequence is not decreasing and is not a regularization of
a psh limit.  It only shows that class convergence alone does not
make individual FS data fibrewise constant.
No finite-fibre average of these stages on that curve is defined.  The moral
is not that these functions are exotic, but that the finite mechanism must
be applied only after the positive-dimensional fibres have been removed,
which is what the birational step does.

\emph{(3) The determinant construction applied to $f$ is the determinant
construction applied to $h$.}
The third temptation is the most interesting, because it very nearly works.
Although $f$ is not finite, the sheaf $f_*\cO_Y$ \emph{is} a coherent
sheaf of $\cO_X$-algebras of generic rank $d=2$, so one may try to form
determinant norms for it directly and bypass the factorization.  On this
surface the computation is immediate: by Step~3 of
\Cref{ex:ch42-cone-stein}, $\sigma_*\cO_Y=\cO_T$, whence
\[
 f_*\cO_Y=h_*\sigma_*\cO_Y=h_*\cO_T .
\]
The algebra one would be using is the algebra of $h$, on the nose.  Its
determinant norm takes as input a section of a line bundle on $T$, not on
$Y$; the functions $\psi_m$ of part~(2) are not pullbacks from $T$,
as their varying values on $E_\sigma$ show.
By contrast $\sigma^*\log|u-w|$ is an explicitly known pullback
and can be transferred to its section data on $T$.
For an arbitrary psh function in a pulled-back class, that transfer
requires the birational reflection/carrier argument; the structure
algebra identity does not prove it.  The apparently one-step
argument therefore contains the two-step argument, with its birational half
concealed inside the identity $f_*\cO_Y=h_*\cO_T$.  This is the assertion
of \Cref{rem:ch42-not-one-step}.

Two further observations complete the picture.  First, the identity
$f_*\cO_Y=h_*\cO_T$ is not special to this surface: it is the definition of
the Stein factorization, $T=\Spec_X(f_*\cO_Y)$, so the coincidence holds
for every proper $f$.  What is special here is only that both sides can be
written down.  Second, the algebra $h_*\cO_T$ is not locally free at the
vertex --- the fibre computation of Step~4 of \Cref{ex:ch42-cone-stein} gives
length three over a point of generic rank two --- so even after reducing to
$h$ one still needs the flattening of \Cref{lem:bj-finite-flatification}
for the constant-rank, arbitrary-base-change determinant construction.
A generic norm on the normal target can still exist.
The proved general descent uses both diagrams in the displayed order;
this proof requirement is not a universal impossibility claim
about other constructions.
\end{example}

\section{Proper support, the extended-real audit, and checkpoints}
\label{sec:ch42-audit}

\subsection{Where extended-real values enter, and where they do not}
\label{sec:ch42-extended-real}

The proof of \Cref{thm:ch42-alteration-descent}, and hence the proof of
\Cref{thm:bj-normal-envelope}, manufactures no psh function and performs no
limit of psh functions: it invokes the two previously established descent
theorems.  Its extended-real bookkeeping is correspondingly short.  The
worked examples do make pointwise calculations, which are audited
separately where they occur.

In the descent proof there are four places where extended-real values must
be kept in view.

\emph{(1) The hypothesis of the envelope property.}  In
\Cref{def:bj-envelope-property} the family $(\varphi_i)_{i\in I}$ is
required to be nonempty and bounded above.  ``Bounded above'' means by a
real constant, not by $+\infty$; each individual $\varphi_i$ may take the
value $-\infty$.  The algebraic common-zero description in
\Cref{prop:ch34-generic-finiteness} applies to FS stages, not to
an arbitrary psh polar locus.  No member of the family is required
to be finite anywhere except at $v_{X,\triv}$, where generic finiteness
forces a real value.  Every application of a descent theorem in this chapter
passes the family on unchanged, so this hypothesis is transported and never
strengthened.

\emph{(2) The regularized supremum.}  The object $(\sup_i\varphi_i)^*$ is
a function with values in $[-\infty,+\infty)$; its value at
$v_{X,\triv}$ is real, being between $\varphi_{i_0}(v_{X,\triv})$ and the
bound.  It is never subtracted from anything in this chapter.

\emph{(3) Pullback along the two factors.}  For a function
$a:X^{\An}\to[-\infty,+\infty)$ the pullbacks $h^*a$ and $g^*h^*a$ are
defined pointwise, with no arithmetic.
For $h$, \Cref{lem:bj-finite-normal-open} gives commutation with
upper-semicontinuous regularization.
No such assertion for an arbitrary function along $g$ is needed:
the birational EP proof uses source EP, carrier reflection and the
divisorial increasing-net criterion, then identifies the regularized
supremum by that criterion.  It does not assume birational analytic
maps to be open.

\emph{(4) The class ledger.}  The identity
\eqref{eq:bj-alteration-class-ledger} is an identity in a finite-dimensional
real vector space and involves no extended reals at all.  This is worth
noting: the assembly proof performs no approximation of classes, so the
$O(m^{-1})$ drift of \Cref{chap:bj-smooth-characteristic} and the varying
classes of \Cref{def:bj-psh} are handled entirely inside the chapters that
are quoted, as \Cref{rem:ch42-real-classes} records.

\begin{warning}{Two operations that this chapter never performs}
\label{warn:ch42-forbidden}
Neither of the following appears above, and each is either undefined or
false.

First, an averaging operator for the alteration.  No finite-fibre
$\operatorname{Av}_f$ is constructed for a nonfinite alteration.  The operator of \Cref{sec:ch41-average-operator}
is constructed from a finite locally free algebra, and the algebra attached
to $f$ is $f_*\cO_Y=h_*\cO_T$, which is the algebra of the finite
factor; the operator built from it consumes Fubini--Study data on $T$, not
on $Y$.  A formula of the shape
$\frac1d\sum_{y\in f^{-1}(x)}e(y/x)\,u(y)$ has no meaning when
$f^{-1}(x)$ is positive-dimensional, and
\Cref{ex:ch42-no-separation}\,(2) exhibits a Fubini--Study function whose
values on such a fibre are not constant.

Second, a difference of descended functions.  The proof of
\Cref{thm:ch42-alteration-descent} uses only implications between envelope
properties and never compares two psh functions by subtraction; in
particular the inequality $u\le P+C$ that carries the tail control in
\Cref{chap:bj-reflection} and \Cref{chap:bj-finite-descent} is used there
and is never converted here into a statement about $u-P$.  Where a carrier
takes the value $-\infty$ the dominated function does too, and the
inequality degenerates to $-\infty\le-\infty+C$; the audits of
\Cref{sec:ch39-audit} and \Cref{sec:ch41-audit} contain the details, and
nothing in the present chapter reopens them.
\end{warning}

\subsection{Values at proper-support semivaluations}
\label{sec:ch42-proper-support}

The envelope property is a statement about functions on the whole of
$X^{\An}$, including the semivaluations whose support is a nonzero prime,
and the chapter must not silently restrict to $X^{\divis}$ or to
$X^{\val}$.  Four checks are relevant.

First, the geometric statements are scheme-theoretic and therefore see every
support stratum.  \Cref{lem:ch42-generic-finite},
\Cref{lem:ch42-stein-universal}, \Cref{lem:bj-alteration-middle},
\Cref{lem:ch42-middle-normalization} and \Cref{lem:ch42-h-degree} are
statements about morphisms of schemes; the analytifications appear only
through $f^{\An}=h^{\An}\circ g^{\An}$, which is a statement about all
points.

Second, the identification at the end of the descent is global.  Step~3 of
\Cref{thm:ch42-alteration-descent} concludes with
$\operatorname{EP}_X(\theta)$, whose conclusion is a membership in
$\PSH_X(\theta)$, and membership is a condition at every point of
$X^{\An}$.  Nothing is inferred from equality on $X^{\divis}$ alone; the
promotion of divisorial equality supplied by
\Cref{prop:ch39-determination-package} is not used in this chapter, and it
would not certify an equality between a psh function and an
otherwise uncertified upper-semicontinuous function.

Third, the surjectivity used is analytic.  Wherever a function is recovered
from its pullback --- inside \Cref{thm:bj-finite-normal-descent}, and in the
verification $u_X=v$ --- the tool is surjectivity of the relevant
analytification, not density of a divisorial locus.  For the finite factor
this is \Cref{lem:bj-finite-normal-open}, which is proved at every analytic
point, including points with nonzero support, and whose proof takes place
inside the residue field of the support prime; the scope statement is
\Cref{rem:ch41-openness-scope}.

Fourth, the worked examples respect this.  In \Cref{ex:ch42-endtoend} the
identity $\sup_c\varphi_c=u_X$ was proved at an arbitrary
$x\in X^{\An}$ by an argument inside $\mathcal H(x)$, with the cases
``$z$ vanishes on the support of $x$'' and ``$a^2$ vanishes on the
support of $x$'' treated explicitly; both are proper-support phenomena
when the support is nonzero.  A vanishing term contributes $-\infty$ and
the other term gives the common maximum on both sides; both sides are
$-\infty$ only when both sections vanish.  In
\Cref{ex:ch42-inseparable-cone}\,(e) the vanishing coefficient
$c_{(1,1)}=0$ contributes $-\infty$ to a maximum and is discarded by the
maximum, not by a convention.

\subsection{Checkpoints}
\label{sec:ch42-checkpoints}

\begin{enumerate}
\item The alteration is smooth, not merely normal.  Smoothness is used
      once, to make \Cref{thm:bj-smooth-envelope} applicable; replacing it
      by normality would make the argument circular.
\item The Stein factorization is formed before anything analytic happens.
      Its middle term is normal, and normality comes from the
      relative-normalization theorem and not from Stein factorization.
\item The birational descent is applied first and the finite descent
      second.  This order is forced by the shape of the factorization, not
      chosen.
\item No degree factor is inserted in the birational step.  The two
      occurrences of $1/d$ are both inside the finite step, and both are
      divisions of real objects.
\item No hypothesis on $d$ is used: not primality, not invertibility in
      $k$, not separability.  \Cref{ex:ch42-inseparable-cone} exercises
      the case $\operatorname{char}k\mid d$ with an inseparable extension.
\item The class $\theta$ is real and arbitrary.  No positivity is
      assumed at any stage, and the statement is vacuous exactly when
      $\PSH_X(\theta)=\varnothing$, by \Cref{lem:ch34-psh-nonempty}.
\item The general descent theorem assumes EP on its normal source.
      In the normal-case proof that source is smooth, so
      \Cref{thm:bj-smooth-envelope} supplies the hypothesis.
      \Cref{thm:bj-final-envelope} is not used, in accordance with
      \Cref{warn:ch42-firewall}.
\item The two corollaries asserting that pullback is bijective are stated
      after the theorem and used nowhere before it.
\end{enumerate}

\section{Failure boundaries}
\label{sec:ch42-failure}

The following boundaries record where the available proof uses the
hypotheses of \Cref{thm:bj-normal-envelope} and its structural choices.
An obstruction to that proof is distinguished from a counterexample
to the conclusion.

\emph{Drop normality of the target.}  The finite-descent proof cannot
be applied unchanged.  On a nodal normalization, $h^{\An}$ is not
open and does not commute with regularization of every bounded
function, as \Cref{ex:ch41-node-not-open} proves.  That example
does not itself give a psh-family counterexample to EP;
\Cref{sec:ch37-node-cusp} makes the same distinction for
branch-asymmetric functions.  The extension to unibranch nonnormal
targets is established in
\Cref{chap:bj-universal-homeomorphisms,chap:bj-subintegral-normalization}.

\emph{Drop normality of the middle term.}  Both descent theorems require
normal source and target, and neither would apply.  Normality of $T$ is
supplied by
\cite[\href{https://stacks.math.columbia.edu/tag/035L}{Tag 035L}]
{stacks-project} through \Cref{prop:appB-stein-normal}, using normality of
$Y$; it is not part of
\cite[\href{https://stacks.math.columbia.edu/tag/03H0}{Tag 03H0}]
{stacks-project}.  \Cref{rem:ch42-normal-middle} lists the four separate
points, in two different chapters, at which it is consumed.

\emph{Drop smoothness of the alteration source.}  Then
\Cref{thm:bj-smooth-envelope} is unavailable and the leftmost hypothesis of
the chain \eqref{eq:bj-alteration-descent-chain} has no supplier.  The
temptation is to observe that a de Jong source is in particular normal and to
invoke the normal case; that is the circle described in
\Cref{rem:ch42-firewall-not-vacuous}, and it proves nothing.  Note that the
descent theorem \Cref{thm:ch42-alteration-descent} is nevertheless stated for
a normal source, because that is all its proof uses; smoothness enters only
where the hypothesis is discharged.

\emph{Drop the Stein factorization.}  Then one is trying to descend along a
morphism that satisfies the hypotheses of neither transport theorem, and
\Cref{ex:ch42-no-separation} shows in three different ways what goes wrong:
there is no common function field for the b-divisor machinery, there are
positive-dimensional fibres on which Fubini--Study functions are not
constant, and the finite algebra one might use is by definition the algebra
of the finite Stein factor.

\emph{Reverse the order of the two descents.}  This needs a
factorization with the finite morphism on top.
Stein factorization does not supply it for the given source.
\Cref{ex:ch42-second-factorization} gives an explicit instance
with a different middle variety.  When such a factorization exists,
normalizing its middle term allows the reverse route, as
\Cref{warn:ch42-order} explains.

\emph{Replace the alteration by a resolution.}  Legitimate in
characteristic zero, where it gives the short proof
\Cref{cor:ch42-char-zero-shortcut}, and unavailable in characteristic $p$
as a uniform all-dimensional input.  The scope of the shortcut is
\Cref{rem:ch42-shortcut-scope}: it removes
\Cref{chap:bj-finite-descent} from the proof of the normal case only.

\emph{Assume every alteration is Galois, or separable.}  Neither is true.
De Jong's theorem permits a generically étale choice over the present
perfect field, but the descent theorem is deliberately stated for an
arbitrary alteration and does not use that choice.
\Cref{ex:ch42-inseparable-cone} exhibits an inseparable finite factor for an
explicit alteration of a singular normal surface.

\emph{Average by the trace, or by a group.}  The trace of a purely
inseparable extension of degree divisible by the characteristic vanishes
identically, and there is no group to average over;
\Cref{ex:ch42-inseparable-cone}\,(d) computes both failures, and
\Cref{rem:ch41-norm-vs-trace} explains why the determinant, being
multiplicative, converts the ultrametric \emph{identity} for products of
Gauss norms into an average, whereas the trace converts the ultrametric
\emph{inequality} for sums into a bound in the wrong direction.

\emph{Insert a degree factor in the birational step.}  Then the classes fail
to match at the middle term, since $g^*$ preserves intersection numbers
while $h^*$ multiplies them by $d$;
\Cref{warn:ch42-no-degree} and
\Cref{lem:ch42-class-ledger}\,\cref{itm:ch42-cl:degree}, with the numbers in
Step~5 of \Cref{ex:ch42-cone-stein}.

\emph{Conflate the Stein factorization with the flattening square.}  They
are different diagrams, built for different purposes, and both occur, one
inside the other;
\Cref{warn:ch42-stein-not-flattening} and the two computations
\Cref{ex:ch42-cone-stein} and \Cref{ex:ch41-flatten-cone}, whose middle
terms are $\mathbf P^2$ and $\mathbf F_2$ respectively.
\Cref{ex:ch42-endtoend} runs both on the same family and keeps the two
copies of $\mathbf F_1$ apart.

\section{Chapter dictionary}
\label{sec:ch42-dictionary}

The correspondences established or used in this chapter are the following.
Each row records a statement together with the hypotheses under which it was
proved.  Rows one to three concern the morphism, rows four to six the
factorization, and rows seven and eight the transport of the envelope
property.

\begin{center}
\begin{tabular}{p{.29\textwidth}p{.31\textwidth}p{.29\textwidth}}
\toprule
geometric datum & algebraic datum & hypothesis\\
\midrule
an alteration $f:Y\to X$ &
proper, surjective, with $K(Y)/K(X)$ finite of degree $d$ &
$X,Y$ integral
(\Cref{def:bj-alteration}, \Cref{lem:ch42-generic-finite})\\
a smooth projective source &
regular over a perfect field, compactified with an SNC boundary; the
alteration may be chosen generically étale &
$k$ algebraically closed, so perfect and without finite extensions
(\Cref{thm:bj-de-jong-input})\\
the degree $d$ &
$[K(Y):K(X)]_{\mathrm s}\cdot[K(Y):K(X)]_{\mathrm i}$, multiplicative in
towers &
dominant generically finite morphisms of integral varieties
(\Cref{def:ch42-alteration-degree}, \Cref{lem:ch42-degree-composite})\\
the Stein middle term $T$ &
$\Spec_X(f_*\cO_Y)$, the integral closure of $\cO_X$ in $K(Y)$ &
$f$ proper dominant generically finite, $Y$ normal integral
(\Cref{lem:ch42-stein-universal}, \Cref{lem:ch42-middle-normalization})\\
normality of $T$ &
relative normalization in a normal source is normal, Tag 035L &
$Y$ normal, $f$ quasi-compact quasi-separated
(\Cref{prop:appB-stein-normal}, \Cref{lem:bj-alteration-middle})\\
the separation $f=h\circ g$ &
$g_*\cO_Y=\cO_T$ with $K(T)=K(Y)$; $h$ finite of degree $d$ &
$X,Y$ normal integral projective, $f$ an alteration
(\Cref{prop:ch42-separation}, \Cref{lem:ch42-h-degree})\\
the class ledger &
$f^*=g^*h^*$, with $g^*$ preserving and $h^*$ multiplying
intersection numbers by $d$ &
integral projective of the same dimension
(\Cref{lem:ch42-class-ledger}, \Cref{warn:ch42-no-degree})\\
$\operatorname{EP}_Y(f^*\theta)\Rightarrow\operatorname{EP}_X(\theta)$ &
birational descent along $g$, then finite descent along $h$ &
one-way as proved; no hypothesis on $d$, on separability, on the
characteristic or on positivity
(\Cref{thm:ch42-alteration-descent})\\
\bottomrule
\end{tabular}
\end{center}

Five separations should be carried forward.

\emph{Alteration versus resolution.}  A resolution is birational; an
alteration is allowed a finite extension of the function field, and that is
the whole reason it exists in every characteristic.  The price is paid
exactly once, in the finite Stein factor, and the degree of that factor is
the degree of the alteration
(\Cref{def:bj-alteration}, \Cref{lem:ch42-h-degree}).  In characteristic
zero the price can be avoided, at the cost of a theorem that is unavailable
in general (\Cref{cor:ch42-char-zero-shortcut}).

\emph{Stein factorization versus flattening square.}  The first is a
triangle factoring a proper morphism into a connected-fibre map
followed by a finite map; the first map is birational in the generically
finite case used here.  The second is a square
built out of a given finite morphism, modifying both source and target in
order to produce a finite \emph{flat} morphism, and it exists by
Raynaud--Gruson flattening.  Both occur in the proof of
\Cref{thm:bj-normal-envelope}, the second inside the first
(\Cref{warn:ch42-stein-not-flattening}).

\emph{Connected fibres versus birationality.}  Stein factorization delivers
$g_*\cO_Y=\cO_T$; birationality of $g$ is deduced from it by computing
the generic fibre.  The two are not the same statement, and in positive
characteristic they come apart: a finite radicial surjective morphism has
connected --- indeed singleton --- fibres and is not birational
when its generic degree exceeds one (\Cref{rem:ch42-connected-not-birational},
\Cref{ex:ch42-inseparable-cone}).

\emph{The degree as a real normalization versus a division in $k$.}  The
number $d$ enters twice, both times inside the finite descent: as the
factor $1/(md)$ normalizing a real-valued maximum, and as the factor
$1/d$ in a numerical class.  Neither is a division in $k$ or in a finite
algebra, so both survive $\operatorname{char}k\mid d$
(\Cref{warn:ch42-no-degree}, \Cref{ex:ch42-inseparable-cone}\,(g)).

\emph{Envelope property as hypothesis versus as theorem.}  In this chapter
the general descent theorem assumes EP on a normal source.
For the normal-case proof, \Cref{thm:bj-smooth-envelope} supplies
that hypothesis on the chosen smooth alteration.
The descent steps then establish EP on the middle term and base.  The two bijectivity statements
\Cref{cor:ch42-birational-bijection} and \Cref{cor:ch42-finite-bijection}
are outputs and are never used as inputs
(\Cref{warn:ch42-firewall}, \Cref{rem:ch42-bijections-are-output}).

\section{Source note}
\label{sec:ch42-source-note}

The architecture of this chapter --- de Jong alteration, Stein
factorization, birational descent on the upper factor and finite descent
on the lower factor --- follows the corresponding section of the supplied
unpublished envelope manuscript.  Its provenance and availability boundary
is recorded in \Cref{sec:appL-envelope}.
This chapter verifies the geometric hypotheses and numerical interfaces
of that assembly explicitly.  No fresh reading of the original supplied
archive or previously published status for its envelope result is claimed.

The precise external attributions are as follows.  The alteration theorem is
de Jong's \cite[Theorem~4.1]{deJong1996}; its final sentence supplies the
generically étale choice over a perfect field, while Remark~4.2 supplies the
finite-ground-field and smoothness refinements.  These are imported through
\Cref{thm:appB-de-jong} and restated with hypotheses displayed as
\Cref{thm:bj-de-jong-input}.  Stein factorization over a
noetherian base is
\cite[\href{https://stacks.math.columbia.edu/tag/03H0}{Tag 03H0}]
{stacks-project}, imported as \Cref{thm:appB-stein-factorization}; normality
of the relative normalization in a normal source is
\cite[\href{https://stacks.math.columbia.edu/tag/035L}{Tag 035L}]
{stacks-project}, imported as \Cref{prop:appB-stein-normal}; finiteness of
the normalization of a Nagata scheme, used in
\Cref{cor:bj-normalization-source-ep}, is
\cite[\href{https://stacks.math.columbia.edu/tag/035S}{Tag 035S}]
{stacks-project}.  These three tags are audited in \Cref{app:stacks}, with
the exact statement, the hypotheses to be checked and the first place of
application; the dependency diagram is \Cref{sec:stacks-dependency-chains}.
No tag is quoted here that is not listed there.  Resolution of singularities
in characteristic zero, used only for the shortcut
\Cref{cor:ch42-char-zero-shortcut}, is \cite{Hironaka1964I} and
\cite{Hironaka1964II}, imported as \Cref{thm:appB-resolution}.  Numerical
pullback, the projection formula and the behaviour of intersection numbers
under a generically finite morphism are those of
\Cref{app:intersection-bdivisors}.

The three mathematical inputs are proved in this book:
\Cref{thm:bj-smooth-envelope} in \Cref{chap:bj-smooth-characteristic},
\Cref{thm:bj-birational-ep-descent} in \Cref{chap:bj-reflection}, and
\Cref{thm:bj-finite-normal-descent} in \Cref{chap:bj-finite-descent}.  The
psh formalism is that of \cite{BJ22}, in the corrected form of
\Cref{def:bj-psh} in which the approximating classes are allowed to vary;
\Cref{ex:ch42-no-separation}\,(2) checks that individually admissible
FS stages can have drifting classes and nonconstant fibre values;
its sequence is not a decreasing regularization and does not prove
that class drift is necessary for a particular psh function.  The envelope property
was formulated in \cite{BJ22} as Boucksom--Jonsson's Conjecture~5.14;
\cite{BJ22b} records partial results available in that literature.  The proof
developed in this book is \Cref{thm:bj-final-envelope}, with historical
context in \Cref{chap:bj-envelope-history}.

Several results cited by the supplied manuscript in neighbouring sections
have no counterpart here and are not cited: Raynaud--Gruson flattening
\cite{RaynaudGruson1971}, the determinant norm functor \cite{FerrandNorm},
the additive-power theorem \cite{RobertsSingh} and the test-ideal
discrepancy bound \cite{Canton} all belong to
\Cref{chap:bj-finite-descent}, \Cref{chap:bj-smooth-characteristic} and
\Cref{chap:bj-subintegral-normalization}, and this chapter consumes them
only through those chapters.

Six aspects of the assembly are isolated here.  \Cref{lem:ch42-generic-finite} proves three general
formulations of generic finiteness equivalent and, under properness, adds
the fourth formulation ``finite over a nonempty open'', because all four
are used below and the literature moves between them without comment.
\Cref{lem:ch42-stein-universal} isolates the universal property of the
middle term and proves it from the adjunction defining the relative
spectrum, which is what makes ``the birational part'' and ``the finite
part'' of an alteration well defined; \Cref{ex:ch42-second-factorization}
shows that without it the phrase would be ambiguous.
\Cref{lem:ch42-middle-normalization} identifies the middle term with an
integral closure, which is the form one computes with.
\Cref{prop:ch42-separation} collects the separation as a single statement
whose items match the hypotheses of the two descent theorems one by one, so
that \Cref{sec:ch42-descent} proceeds by quotation and the verification is
visible in one place.  \Cref{lem:ch42-class-ledger} and
\Cref{warn:ch42-no-degree} write out the numerical bookkeeping, including
the assertion that no degree factor belongs in the birational step; this is
the point at which a plausible argument most easily goes wrong.  Finally the
worked computations \Cref{ex:ch42-frobenius-alteration},
\Cref{ex:ch42-degree-matters}, \Cref{ex:ch42-not-a-pullback},
\Cref{ex:ch42-cone-stein}, \Cref{ex:ch42-stein-degenerate},
\Cref{ex:ch42-second-factorization}, \Cref{ex:ch42-endtoend},
\Cref{ex:ch42-inseparable-cone} and \Cref{ex:ch42-no-separation}
are worked checks, included because the two claims this chapter exists to establish
--- that the two mechanisms can be separated, and that the separation costs
nothing numerically --- are exactly the claims that a summary can appear to
establish without doing so.  A reader who has computed the Stein
factorization of an explicit alteration of $\mathbf P(1,1,2)$, followed a
family of psh functions through all five varieties, and watched the trace
vanish in characteristic two while the determinant did not, will not confuse
the two diagrams again.
